\input{preamble}

% OK, start here.
%
\begin{document}

\title{Théorie de l'intersection}


\maketitle

\phantomsection
\label{section-phantom}

\tableofcontents


\section{Introduction}
\label{section-introduction}

\noindent
Dans ce chapitre, nous construisons le produit d'intersection sur les groupes de Chow
modulo l'équivalence rationnelle d'une variété projective non singulière sur un
corps algébriquement clos. Nos outils sont la formule des Tor de Serre
(voir \cite[chapitre V]{Serre_algebre_locale}), la réduction à la diagonale
et le lemme de déplacement.

\medskip\noindent
Nous rappelons d'abord les cycles et la construction de l'image directe propre et de
l'image inverse plate des cycles. Ensuite, nous introduisons l'équivalence rationnelle des cycles,
qui nous donne les groupes de Chow $\CH_*(X)$. L'image directe propre et l'image inverse plate
se factorisent par l'équivalence rationnelle et donnent des opérations sur les groupes de Chow.
Cela occupe les sections
\ref{section-cycles},
\ref{section-cycle-of-closed},
\ref{section-cycle-of-coherent-sheaf},
\ref{section-pushforward},
\ref{section-flat-pullback},
\ref{section-rational-equivalence},
\ref{section-alternative},
\ref{section-pushforward-and-rational-equivalence}, et
\ref{section-flat-pullback-and-rational-equivalence}.
Pour les démonstrations, nous renvoyons principalement au chapitre sur l'homologie de Chow,
où ces résultats ont été démontrés dans le cadre des
schémas localement de type fini sur une base noethérienne universellement caténaire ; voir
Homologie de Chow, section \ref{chow-section-setup} et suivantes.

\medskip\noindent
Puisque nous travaillons sur une variété projective non singulière $X$, toute composante irréductible
de l'intersection $V \cap W$ de deux sous-variétés fermées irréductibles
est de dimension au moins $\dim(V) + \dim(W) - \dim(X)$. Nous disons que $V$ et $W$
se coupent proprement si l'égalité vaut pour toute composante irréductible $Z$.
Dans ce cas, nous définissons la multiplicité d'intersection
$e_Z = e(X, V \cdot W, Z)$ par la formule
$$
e_Z = \sum\nolimits_i
(-1)^i
\text{length}_{\mathcal{O}_{X, Z}}
\text{Tor}_i^{\mathcal{O}_{X, Z}}(\mathcal{O}_{W, Z}, \mathcal{O}_{V, Z})
$$
Nous devons faire un peu d'algèbre commutative pour montrer que ces
multiplicités d'intersection concordent avec l'intuition dans les cas simples,
c'est-à-dire que, parfois,
$$
e_Z = \text{length}_{\mathcal{O}_{X, Z}} \mathcal{O}_{V \cap W, Z},
$$
autrement dit, seul $\text{Tor}_0$ contribue. Cela se produit lorsque
$V$ et $W$ sont de Cohen-Macaulay au point générique de $Z$, ou lorsque
$W$ est défini par une suite régulière dans $\mathcal{O}_{X, Z}$ qui
définit aussi une suite régulière sur $\mathcal{O}_{V, Z}$. Toutefois,
l'exemple \ref{example-naive-multiplicity-wrong} montre que les Tor supérieurs
sont nécessaires en général. De plus, il existe un lien
avec la multiplicité de Samuel. Ces questions sont étudiées dans les
sections
\ref{section-intersect-properly},
\ref{section-tor-formula},
\ref{section-multiplicities},
\ref{section-computing-intersection-multiplicities}, et
\ref{section-intersection-product}.

\medskip\noindent
La réduction à la diagonale affirme que nous pouvons couper
$V$ et $W$ en coupant $V \times W$ avec la diagonale dans $X \times X$.
Cet énoncé anodin, évident au niveau des intersections
schématiques, ramène l'intersection d'une paire quelconque
de sous-schémas fermés au cas où l'un des deux est localement défini
par une suite régulière. À la suite de Serre, nous l'utilisons pour obtenir la positivité
des multiplicités d'intersection. De plus, la réduction à la diagonale
conduit à l'additivité des multiplicités d'intersection, à l'associativité et à une
formule de projection. On trouvera cela dans les sections
\ref{section-exterior-product},
\ref{section-reduction-diagonal},
\ref{section-associative},
\ref{section-flat-pullback-and-intersection-products}, et
\ref{section-projection-formula-flat}.

\medskip\noindent
Enfin, nous abordons les lemmes de déplacement et leurs applications. Le lemme de
déplacement comporte deux parties. La première affirme qu'étant données des sous-variétés fermées
$$
Z \subset X \subset \mathbf{P}^N
$$
avec $X$ non singulière, nous pouvons trouver une sous-variété $C \subset \mathbf{P}^N$
qui coupe $X$ proprement et telle que
$$
C \cdot X = [Z] + \sum m_j [Z_j]
$$
et telle que les autres composantes $Z_j$ soient « plus générales » que $Z$.
La seconde partie affirme que l'on peut déplacer $C \subset \mathbf{P}^N$, le long
d'une courbe rationnelle, vers une sous-variété en position générale par rapport à
toute liste donnée de sous-variétés. Ensemble, ces résultats impliquent qu'il suffit
de définir le produit d'intersection des cycles sur $X$ qui se coupent
proprement, ce qui a été fait ci-dessus. Bien entendu, cela ne conduit à un produit d'intersection
sur $\CH_*(X)$ que si l'on montre, comme nous le faisons dans le texte, que ces produits
passent au quotient par l'équivalence rationnelle. Cette question et quelques applications sont étudiées
dans les sections
\ref{section-projection},
\ref{section-moving-lemma},
\ref{section-intersections-and-rational-equivalence},
\ref{section-chow-rings},
\ref{section-general-pullback}, et
\ref{section-pullback-cycles}.



\section{Conventions}
\label{section-conventions}

\noindent
Nous fixons un corps de base algébriquement clos $\mathbf{C}$, de caractéristique
quelconque. Tous les schémas et toutes les variétés sont définis sur $\mathbf{C}$, et tous les
morphismes sont des $\mathbf{C}$-morphismes. Une variété $X$ est
{\it non singulière} si $X$ est un schéma régulier (voir
Propriétés, définition \ref{properties-definition-regular}).
Dans notre cas, cela signifie que le morphisme $X \to \Spec(\mathbf{C})$
est lisse (voir
Variétés, lemme \ref{varieties-lemma-geometrically-regular-smooth}).


\section{Cycles}
\label{section-cycles}

\noindent
Soit $X$ une variété. Une {\it sous-variété fermée} de $X$ est un sous-schéma
fermé intègre $Z \subset X$. Un {\it $k$-cycle} sur $X$ est une somme formelle
finie $\sum n_i [Z_i]$, où chaque $Z_i$ est une sous-variété fermée
de dimension $k$. Chaque fois que nous employons la notation $\alpha = \sum n_i[Z_i]$
pour un $k$-cycle, nous supposons toujours que les sous-variétés $Z_i$ sont deux à deux
distinctes et que $n_i \not = 0$ pour tout $i$. Dans ce cas, le
{\it support} de $\alpha$ est le sous-ensemble fermé
$$
\text{Supp}(\alpha) = \bigcup Z_i \subset X
$$
de dimension $k$. Le groupe des $k$-cycles est noté $Z_k(X)$.
Voir Homologie de Chow, section \ref{chow-section-cycles}.


\section{Cycle associé à un sous-schéma fermé}
\label{section-cycle-of-closed}

\noindent
Supposons que $X$ soit une variété et que $Z \subset X$ soit un sous-schéma fermé
avec $\dim(Z) \leq k$. Soient $Z_i$ les composantes irréductibles de $Z$ de
dimension $k$, et soit $n_i$ la {\it multiplicité de $Z_i$ dans $Z$},
définie par
$$
n_i = \text{length}_{\mathcal{O}_{X, Z_i}} \mathcal{O}_{Z, Z_i}
$$
où $\mathcal{O}_{X, Z_i}$, resp.\ $\mathcal{O}_{Z, Z_i}$, est l'anneau
local de $X$, resp.\ de $Z$, au point générique de $Z_i$.
Nous définissons le $k$-cycle associé à $Z$ comme le $k$-cycle
$$
[Z]_k = \sum n_i [Z_i].
$$
Voir Homologie de Chow, section \ref{chow-section-cycle-of-closed-subscheme}.


\section{Cycle associé à un faisceau cohérent}
\label{section-cycle-of-coherent-sheaf}

\noindent
Supposons que $X$ soit une variété et que
$\mathcal{F}$ soit un $\mathcal{O}_X$-module cohérent tel que
$\dim(\text{Supp}(\mathcal{F})) \leq k$.
Soient $Z_i$ les composantes irréductibles de $\text{Supp}(\mathcal{F})$
de dimension $k$, et soit $n_i$ la
{\it multiplicité de $Z_i$ dans $\mathcal{F}$}, définie par
$$
n_i = \text{length}_{\mathcal{O}_{X, Z_i}} \mathcal{F}_{\xi_i}
$$
où $\mathcal{O}_{X, Z_i}$ est l'anneau
local de $X$ au point générique $\xi_i$ de $Z_i$,
et où $\mathcal{F}_{\xi_i}$ est le germe de $\mathcal{F}$ en ce point.
Nous définissons le $k$-cycle associé à $\mathcal{F}$ comme le $k$-cycle
$$
[\mathcal{F}]_k = \sum n_i [Z_i].
$$
Voir Homologie de Chow, section \ref{chow-section-cycle-of-coherent-sheaf}.
Notons que, si $Z \subset X$ est un sous-schéma fermé avec $\dim(Z) \leq k$, alors
$[Z]_k = [\mathcal{O}_Z]_k$ par définition.


\section{Image directe propre}
\label{section-pushforward}

\noindent
Supposons que $f : X \to Y$ soit un morphisme propre de variétés.
Soit $Z \subset X$ une sous-variété fermée de dimension $k$.
Nous définissons $f_*[Z]$ comme $0$ si $\dim(f(Z)) < k$,
et comme $d \cdot [f(Z)]$ si $\dim(f(Z)) = k$, où
$$
d = [\mathbf{C}(Z) : \mathbf{C}(f(Z))] = \deg(Z/f(Z))
$$
est le degré du morphisme dominant $Z \to f(Z)$ ; voir
Morphismes, définition \ref{morphisms-definition-degree}.
Soit $\alpha = \sum n_i [Z_i]$ un $k$-cycle sur $X$. L'
{\it image directe} de $\alpha$ est la somme $f_* \alpha = \sum n_i f_*[Z_i]$,
où chaque $f_*[Z_i]$ est défini comme ci-dessus. Cela définit un homomorphisme
$$
f_* : Z_k(X) \longrightarrow Z_k(Y)
$$
Voir Homologie de Chow, section \ref{chow-section-proper-pushforward}.

\begin{lemma}
\label{lemma-push-coherent}
\begin{reference}
Voir \cite[chapitre V]{Serre_algebre_locale}.
\end{reference}
Supposons que $f : X \to Y$ soit un morphisme propre de variétés.
Soit $\mathcal{F}$ un faisceau cohérent tel que
$\dim(\text{Supp}(\mathcal{F})) \leq k$ ; alors
$f_*[\mathcal{F}]_k = [f_*\mathcal{F}]_k$. En particulier, si
$Z \subset X$ est un sous-schéma fermé de dimension $\leq k$, alors
$f_*[Z]_k = [f_*\mathcal{O}_Z]_k$.
\end{lemma}

\begin{proof}
Voir Homologie de Chow, lemme \ref{chow-lemma-cycle-push-sheaf}.
\end{proof}

\begin{lemma}
\label{lemma-compose-pushforward}
Soient $f : X \to Y$ et $g : Y \to Z$ des morphismes propres de
variétés. Alors $g_* \circ f_* = (g \circ f)_*$ comme applications $Z_k(X) \to Z_k(Z)$.
\end{lemma}

\begin{proof}
Cas particulier de Homologie de Chow, lemme \ref{chow-lemma-compose-pushforward}.
\end{proof}



\section{Image inverse plate}
\label{section-flat-pullback}

\noindent
Supposons que $f : X \to Y$ soit un morphisme plat de variétés.
D'après Morphismes, lemme
\ref{morphisms-lemma-dimension-fibre-at-a-point-additive} 
toute fibre de $f$ est de dimension $r = \dim(X) - \dim(Y)$\footnote{Réciproquement,
si $f : X \to Y$ est un morphisme dominant de variétés, si
$X$ est de Cohen-Macaulay, si $Y$ est non singulière et si toutes les fibres ont
la même dimension $r$, alors $f$ est plat. Cela résulte de
Algèbre, lemme \ref{algebra-lemma-CM-over-regular-flat}, et de
Variétés, lemme \ref{varieties-lemma-dimension-fibres-locally-algebraic},
qui montre que $\dim(X) = \dim(Y) + r$.}.
Soit $Z \subset Y$ une sous-variété fermée de dimension $k$. Nous définissons
$f^*[Z]$ comme le $(k + r)$-cycle associé à l'image inverse
schématique : $f^*[Z] = [f^{-1}(Z)]_{k + r}$. Soit
$\alpha = \sum n_i [Z_i]$ un $k$-cycle sur $Y$. L'{\it image inverse} de
$\alpha$ est la somme $f^* \alpha = \sum n_i f^*[Z_i]$, où chaque $f^*[Z_i]$
est défini comme ci-dessus. Cela définit un homomorphisme
$$
f^* : Z_k(Y) \longrightarrow Z_{k + r}(X)
$$
Voir Homologie de Chow, section \ref{chow-section-flat-pullback}.

\begin{lemma}
\label{lemma-pullback}
Soit $f : X \to Y$ un morphisme plat de variétés. Posons $r = \dim(X) - \dim(Y)$.
Alors $f^*[\mathcal{F}]_k = [f^*\mathcal{F}]_{k + r}$
si $\mathcal{F}$ est un faisceau cohérent sur $Y$ et si la dimension du
support de $\mathcal{F}$ est au plus $k$.
\end{lemma}

\begin{proof}
Voir Homologie de Chow, lemme \ref{chow-lemma-pullback-coherent}.
\end{proof}

\begin{lemma}
\label{lemma-compose-flat-pullback}
Soient $f : X \to Y$ et $g : Y \to Z$ des morphismes plats de
variétés. Alors $g \circ f$ est plat et $f^* \circ g^* = (g \circ f)^*$
comme applications $Z_k(Z) \to Z_{k + \dim(X) - \dim(Z)}(X)$.
\end{lemma}

\begin{proof}
Cas particulier de Homologie de Chow, lemme \ref{chow-lemma-compose-flat-pullback}.
\end{proof}


\section{Équivalence rationnelle}
\label{section-rational-equivalence}

\noindent
Nous allons définir l'équivalence rationnelle d'une manière qui peut, au premier
abord, sembler différente de celle qui vous est familière ou de celle
de \cite[chapitre I]{F} ou de
Homologie de Chow, section \ref{chow-section-rational-equivalence}.
Cependant, nous montrerons dans la section \ref{section-alternative} que
les deux notions coïncident.

\medskip\noindent
Soit $X$ une variété. Soit $W \subset X \times \mathbf{P}^1$
une sous-variété fermée de dimension $k + 1$. Soient $a, b$ deux points fermés distincts
de $\mathbf{P}^1$. Supposons que $X \times a$, $X \times b$ et $W$
se coupent proprement :
$$
\dim (W \cap X \times a) \leq k,\quad
\dim (W \cap X \times b) \leq k.
$$
C'est le cas dès que $W \to \mathbf{P}^1$ est dominant, ou si $W$ est
contenu dans une fibre de la projection au-dessus d'un point fermé distinct de
$a$ et de $b$ (cas sans intérêt que nous écarterons). Dans cette
situation, la fibre schématique $W_a$ du morphisme
$W \to \mathbf{P}^1$ est égale à l'intersection schématique
$W \cap X \times a$ dans $X \times \mathbf{P}^1$. En identifiant $X \times a$
et $X \times b$ à $X$, nous pouvons considérer les fibres $W_a$ et $W_b$
comme des sous-schémas fermés de $X$ de dimension $\leq k$\footnote{Nous considérerons parfois
$W_a$ comme un sous-schéma fermé de $X \times \mathbf{P}^1$, et parfois
comme un sous-schéma fermé de $X$. Le contexte indiquera toujours clairement quel
point de vue est adopté.}. Un exemple fondamental d'
équivalence rationnelle est
$$
[W_a]_k \sim_{rat} [W_b]_k
$$
Les cycles $[W_a]_k$ et $[W_b]_k$ se calculent facilement en pratique
(une fois $W$ donné), car ils s'obtiennent par intersection propre avec
un diviseur de Cartier (nous le verrons dans la
section \ref{section-intersection-product}).
Comme le groupe des automorphismes de $\mathbf{P}^1$ est $2$-transitif, nous pouvons
envoyer la paire de points fermés $a, b$ sur toute paire de notre choix. Un choix
traditionnel consiste à prendre $a = 0$ et $b = \infty$.

\medskip\noindent
Plus généralement, soit $\alpha = \sum n_i [W_i]$ un $(k + 1)$-cycle sur
$X \times \mathbf{P}^1$. Soient $a_i, b_i$ des paires de points fermés distincts de
$\mathbf{P}^1$. Supposons que $X \times a_i$, $X \times b_i$ et $W_i$ se coupent
proprement ; autrement dit, chaque triplet $W_i, a_i, b_i$ vérifie la condition
étudiée ci-dessus. Un {\it cycle rationnellement équivalent à zéro} est tout cycle
de la forme
$$
\sum n_i([W_{i, a_i}]_k - [W_{i, b_i}]_k).
$$
Il s'agit bien d'un $k$-cycle. L'ensemble des $k$-cycles rationnellement
équivalents à zéro est un sous-groupe additif du groupe des $k$-cycles.
Nous disons que deux $k$-cycles sont {\it rationnellement équivalents}, et nous écrivons
$\alpha \sim_{rat} \alpha'$, si $\alpha - \alpha'$ est un cycle rationnellement
équivalent à zéro.

\medskip\noindent
Nous définissons
$$
\CH_k(X) = Z_k(X)/ \sim_{rat}
$$
comme le {\it groupe de Chow des $k$-cycles sur $X$}. Nous verrons dans le
lemme \ref{lemma-rational-equivalence}
que cela coïncide avec le groupe de Chow défini dans
Homologie de Chow, définition \ref{chow-definition-rational-equivalence}.


\section{Équivalence rationnelle et fonctions rationnelles}
\label{section-alternative}

\noindent
Soit $X$ une variété. Soit $W \subset X$ une sous-variété
de dimension $k + 1$. Soit $f \in \mathbf{C}(W)^*$ une fonction rationnelle non nulle
sur $W$. Pour toute sous-variété $Z \subset W$ de dimension $k$,
on peut définir l'ordre d'annulation $\text{ord}_{W, Z}(f)$ de $f$ le long de
$Z$. Si $f$ appartient à l'anneau local $\mathcal{O}_{W, Z}$,
on a alors
$$
\text{ord}_{W, Z}(f) =
\text{length}_{\mathcal{O}_{X, Z}} \mathcal{O}_{W, Z}/f\mathcal{O}_{W, Z}
$$
où $\mathcal{O}_{X, Z}$, resp.\ $\mathcal{O}_{W, Z}$, est l'anneau
local de $X$, resp.\ de $W$, au point générique de $Z$. En général, on
étend la définition par multiplicativité. Le {\it diviseur principal
associé à $f$} est
$$
\text{div}_W(f) = \sum \text{ord}_{W, Z}(f)[Z]
$$
dans $Z_k(W)$. Comme $W \subset X$ est une sous-variété fermée, nous pouvons considérer
$\text{div}_W(f)$ comme un cycle sur $X$.
Voir Homologie de Chow, section \ref{chow-section-principal-divisors}.

\begin{lemma}
\label{lemma-rational-equivalence}
Soit $X$ une variété. Soit $W \subset X$ une sous-variété
de dimension $k + 1$. Soit $f \in \mathbf{C}(W)^*$ une fonction rationnelle non nulle
sur $W$. Alors $\text{div}_W(f)$ est rationnellement équivalent à zéro sur
$X$. Réciproquement, ces diviseurs principaux engendrent le groupe abélien des
cycles rationnellement équivalents à zéro sur $X$.
\end{lemma}

\begin{proof}
La première assertion résulte de
Homologie de Chow, lemme \ref{chow-lemma-rational-function}.
Plus précisément, soit $W' \subset X \times \mathbf{P}^1$ l'adhérence
du graphe de $f$. Alors $\text{div}_W(f) = [W'_0]_k - [W'_\infty]$
dans $Z_k(W) \subset Z_k(X)$ ; voir la partie (6) de
Homologie de Chow, lemme \ref{chow-lemma-rational-function}.

\medskip\noindent
Pour la seconde assertion, soit $W' \subset X \times \mathbf{P}^1$ une sous-variété
fermée de dimension $k + 1$ qui domine $\mathbf{P}^1$.
Nous allons montrer que $[W'_0]_k - [W'_\infty]_k$ est un diviseur principal,
ce qui achèvera la démonstration. Soit $W \subset X$ l'image de $W'$
par la projection sur $X$. Alors $W \subset X$ est une sous-variété fermée
et $W' \to W$ est propre et dominant, à fibres de dimension $0$
ou $1$. Si $\dim(W) = k$, alors $W' = W \times \mathbf{P}^1$ et nous
voyons que $[W'_0]_k - [W'_\infty]_k = [W] - [W] = 0$.
Si $\dim(W) = k + 1$, alors
$W' \to W$ est génériquement fini\footnote{Si $W' \to W$ est birationnel,
le résultat résulte de
Homologie de Chow, lemme \ref{chow-lemma-rational-function}.
Notre tâche est de montrer que, même si $W' \to W$
est de degré $> 1$, l'équivalence rationnelle fondamentale
$[W'_0]_k \sim_{rat} [W'_\infty]_k$ provient d'un diviseur principal
sur une sous-variété de $X$.}. Notons $f$ la projection $W' \to \mathbf{P}^1$,
considérée comme un élément de $\mathbf{C}(W')^*$. Soit
$g = \text{Nm}(f) \in \mathbf{C}(W)^*$ la norme. D'après
Homologie de Chow, lemme \ref{chow-lemma-proper-pushforward-alteration},
nous avons
$$
\text{div}_W(g) = \text{pr}_{X, *}\text{div}_{W'}(f)
$$
Comme $\text{div}_{W'}(f) = [W'_0]_k - [W'_\infty]_k$
d'après Homologie de Chow, lemme \ref{chow-lemma-rational-function},
la démonstration est achevée.
\end{proof}


\section{Image directe propre et équivalence rationnelle}
\label{section-pushforward-and-rational-equivalence}

\noindent
Supposons que $f : X \to Y$ soit un morphisme propre de variétés.
Soit $\alpha \sim_{rat} 0$ un $k$-cycle sur
$X$ rationnellement équivalent à $0$. Alors l'{image directe}
de $\alpha$ est rationnellement équivalente à zéro :
$f_* \alpha \sim_{rat} 0$. Voir le chapitre I de \cite{F} ou
Homologie de Chow, lemme \ref{chow-lemma-proper-pushforward-rational-equivalence}.

\medskip\noindent
Nous obtenons donc un diagramme commutatif
$$
\xymatrix{
Z_k(X) \ar[r] \ar[d]_{f_*} & \CH_k(X) \ar[d]^{f_*} \\
Z_k(Y) \ar[r] & \CH_k(Y)
}
$$
de groupes de $k$-cycles.


\section{Image inverse plate et équivalence rationnelle}
\label{section-flat-pullback-and-rational-equivalence}

\noindent
Supposons que $f : X \to Y$ soit un morphisme plat de variétés.
Posons $r = \dim(X) - \dim(Y)$.
Soit $\alpha \sim_{rat} 0$ un $k$-cycle sur
$Y$ rationnellement équivalent à $0$. Alors l'image inverse
de $\alpha$ est rationnellement équivalente à zéro :
$f^* \alpha \sim_{rat} 0$. Voir le chapitre I de \cite{F} ou
Homologie de Chow, lemme \ref{chow-lemma-flat-pullback-rational-equivalence}.

\medskip\noindent
Nous obtenons donc un diagramme commutatif
$$
\xymatrix{
Z_{k + r}(X) \ar[r] & \CH_{k + r}(X) \\
Z_k(Y) \ar[r] \ar[u]^{f^*} & \CH_k(Y) \ar[u]_{f^*}
}
$$
de groupes de $k$-cycles.


\section{La suite exacte courte associée à un ouvert}
\label{section-ses}

\noindent
Soit $X$ une variété et soit $U \subset X$ une sous-variété ouverte.
Soit $X \setminus U = \bigcup Z_i$ sa décomposition en composantes
irréductibles\footnote{Comme, dans ce chapitre, nous ne considérons que les groupes de Chow
de variétés, nous ne pouvons pas prendre $Z_k(X \setminus U)$
et $\CH_k(X \setminus U)$ ; d'où l'approche qui emploie les variétés $Z_i$.}.
Alors, pour tout $k \geq 0$, il existe un diagramme commutatif
$$
\xymatrix{
\bigoplus Z_k(Z_i) \ar[r] \ar[d] &
Z_k(X) \ar[r] \ar[d] &
Z_k(U) \ar[d] \ar[r] &
0 \\
\bigoplus \CH_k(Z_i) \ar[r] &
\CH_k(X) \ar[r] &
\CH_k(U) \ar[r] &
0
}
$$
à lignes exactes. Les flèches verticales sont ici les applications canoniques de passage au quotient.
Les flèches horizontales de gauche sont données par l'image directe propre le long des
immersions fermées $Z_i \to X$. Les flèches horizontales de droite sont données par l'image
inverse plate le long de l'immersion ouverte $j : U \to X$. Comme nous avons vu que
ces applications se factorisent par l'équivalence rationnelle, les carrés
sont commutatifs. La ligne du haut est exacte simplement parce que toute sous-variété
de $X$ est soit contenue dans un des $Z_i$, soit a une intersection irréductible
avec $U$. La ligne du bas est exacte parce que tout diviseur principal
$\text{div}_W(f)$ sur $U$ est la restriction d'un diviseur principal sur $X$.
Plus précisément, si $W \subset U$ est une sous-variété fermée de dimension $(k + 1)$
et si $f \in \mathbf{C}(W)^*$, notons $\overline{W}$ l'adhérence de $W$
dans $X$. Alors $W \subset \overline{W}$ est une immersion ouverte, donc
$\mathbf{C}(W) = \mathbf{C}(\overline{W})$, et nous pouvons considérer $f$
comme une fonction rationnelle non constante sur $\overline{W}$. Alors, clairement,
$$
j^*\text{div}_{\overline{W}}(f) = \text{div}_W(f)
$$
dans $Z_k(X)$. L'exactitude de la ligne inférieure s'en déduit aisément.
Pour les détails, voir Homologie de Chow, lemme \ref{chow-lemma-restrict-to-open}.


\section{Intersections propres}
\label{section-intersect-properly}

\noindent
Commençons par quelques lemmes donnant des estimations de dimension.

\begin{lemma}
\label{lemma-dimension-product-varieties}
Soient $X$ et $Y$ des variétés. Alors $X \times Y$ est une variété et
$\dim(X \times Y) = \dim(X) + \dim(Y)$.
\end{lemma}

\begin{proof}
Le schéma $X \times Y = X \times_{\Spec(\mathbf{C})} Y$ est une variété d'après
Variétés, lemme \ref{varieties-lemma-product-varieties}.
L'assertion sur la dimension est
Variétés, lemme \ref{varieties-lemma-dimension-product-locally-algebraic}.
\end{proof}

\noindent
Rappelons qu'une immersion régulière $i : X \to Y$ de schémas
est une immersion fermée dont le
faisceau d'idéaux correspondant est localement engendré par une suite régulière ; voir
Diviseurs, section \ref{divisors-section-regular-immersions}.
De plus, le faisceau conormal $\mathcal{C}_{X/Y}$ est localement libre de type fini, de
rang égal à la longueur de la suite régulière. Nous dirons que $i$ est une
{\it immersion régulière de codimension $c$}
si $\mathcal{C}_{X/Y}$ est localement libre de rang $c$.

\medskip\noindent
Plus généralement, rappelons
(Morphismes, compléments, section \ref{more-morphisms-section-lci})
que $f : X \to Y$ est un morphisme d'intersection complète locale
si l'on peut recouvrir $X$ par des ouverts $U$ tels que $f|_U$ se factorise
comme suit :
$$
\xymatrix{
U \ar[rr]_i \ar[rd] & & \mathbf{A}^n_Y \ar[ld] \\
& Y
}
$$
où $i$ est une immersion Koszul-régulière (si $Y$ est localement noethérien,
cela revient à demander que $i$ soit une immersion régulière ; voir
Diviseurs, lemme \ref{divisors-lemma-regular-immersion-noetherian}).
Nous dirons que $f$ est un {\it morphisme d'intersection complète locale
de dimension relative $r$} si, pour toute factorisation comme ci-dessus,
l'immersion fermée $i$ a un faisceau conormal de rang $n - r$ (autrement
dit, si $i$ est une immersion Koszul-régulière de codimension $n - r$,
ce qui, dans le cas noethérien, signifie simplement qu'elle est une immersion régulière de
codimension $n - r$).

\begin{lemma}
\label{lemma-pullback-by-regular-immersion}
Soit $f : X \to Y$ un morphisme de variétés.
\begin{enumerate}
\item Si $Z \subset Y$ est une sous-variété de dimension $d$ et si $f$ est une
immersion régulière de codimension $c$, alors toute composante irréductible
de $f^{-1}(Z)$ est de dimension $\geq d - c$.
\item Si $Z \subset Y$ est une sous-variété de dimension $d$ et si
$f$ est un morphisme d'intersection complète locale de dimension relative $r$,
alors toute composante irréductible de $f^{-1}(Z)$ est de dimension $\geq d + r$.
\end{enumerate}
\end{lemma}

\begin{proof}
Démonstration de (1). Nous pouvons travailler localement ; nous pouvons donc supposer que
$Y = \Spec(A)$ et $X = V(f_1, \ldots, f_c)$, où $f_1, \ldots, f_c$
est une suite régulière dans $A$. Si $Z = \Spec(A/\mathfrak p)$, alors
nous voyons que $f^{-1}(Z) = \Spec(A/\mathfrak p + (f_1, \ldots, f_c))$.
Si $V$ est une composante irréductible de $f^{-1}(Z)$, nous pouvons
choisir un point fermé $v \in V$ qui n'appartient à aucune autre composante irréductible
de $f^{-1}(Z)$. Alors
$$
\dim(Z) = \dim \mathcal{O}_{Z, v}
\quad\text{et}\quad
\dim(V) = \dim \mathcal{O}_{V, v} = \dim \mathcal{O}_{Z, v}/(f_1, \ldots, f_c)
$$
La première égalité résulte par exemple de
Algèbre, lemme \ref{algebra-lemma-dimension-prime-polynomial-ring},
et la seconde de notre choix du point fermé.
Le résultat découle alors du fait que le quotient par un élément
de l'idéal maximal diminue la dimension d'au plus $1$ ; voir
Algèbre, lemme \ref{algebra-lemma-one-equation}.

\medskip\noindent
Démonstration de (2). Choisissons une factorisation comme dans la définition d'une
intersection complète locale et appliquons (1). Certains détails sont omis.
\end{proof}

\begin{lemma}
\label{lemma-diagonal-regular-immersion}
Soit $X$ une variété non singulière. Alors la diagonale
$\Delta : X \to X \times X$ est une immersion régulière de codimension $\dim(X)$.
\end{lemma}

\begin{proof}
En fait, toute immersion fermée entre variétés projectives non singulières
est une immersion régulière ; voir Diviseurs,
lemme \ref{divisors-lemma-immersion-smooth-into-smooth-regular-immersion}.
\end{proof}

\noindent
Le lemme suivant illustre le fonctionnement de la réduction à la diagonale.

\begin{lemma}
\label{lemma-intersect-in-smooth}
Soit $X$ une variété non singulière et soient $W,V \subset X$
des sous-variétés fermées telles que $\dim(W) = s$ et $\dim(V) = r$. Alors toute
composante irréductible $Z$ de $V \cap W$ est de dimension $\geq r + s - \dim(X)$.
\end{lemma}

\begin{proof}
Comme $V \cap W = \Delta^{-1}(V \times W)$ (au sens schématique),
nous concluons par les lemmes \ref{lemma-diagonal-regular-immersion} et
\ref{lemma-pullback-by-regular-immersion}.
\end{proof}

\noindent
Ce lemme suggère la définition suivante.

\begin{definition}
\label{definition-proper-intersection}
Soit $X$ une variété non singulière.
\begin{enumerate}
\item {\emergencystretch=1em Soient $W,V \subset X$ des sous-variétés fermées telles que
$\dim(W) = s$ et $\dim(V) = r$. Nous disons que $W$ et $V$
{\it se coupent proprement} si $\dim(V \cap W) \leq r + s - \dim(X)$.\par}
\item Soient $\alpha = \sum n_i [W_i]$ un $s$-cycle
et $\beta = \sum_j m_j [V_j]$ un $r$-cycle sur $X$. Nous disons
que $\alpha$ et $\beta$ {\it se coupent proprement} si
$W_i$ et $V_j$ se coupent proprement quels que soient $i$ et $j$.
\end{enumerate}
\end{definition}


\section{Multiplicités d'intersection au moyen de la formule des Tor}
\label{section-tor-formula}

\noindent
Un fait fondamental que nous utiliserons souvent est que, pour des faisceaux de
modules $\mathcal{F}$, $\mathcal{G}$ sur un espace annelé $(X, \mathcal{O}_X)$
et un point $x \in X$, on a
$$
\text{Tor}_p^{\mathcal{O}_X}(\mathcal{F}, \mathcal{G})_x =
\text{Tor}_p^{\mathcal{O}_{X, x}}(\mathcal{F}_x, \mathcal{G}_x)
$$
comme $\mathcal{O}_{X, x}$-modules. Cela se voit de plusieurs façons
à partir de notre construction des produits tensoriels dérivés dans
Cohomologie, section \ref{cohomology-section-flat}; cela résulte par exemple de
Cohomologie, lemme \ref{cohomology-lemma-check-K-flat-stalks}.
De plus, si $X$ est un schéma et si $\mathcal{F}$ et $\mathcal{G}$
sont quasi-cohérents, alors les modules
$\text{Tor}_p^{\mathcal{O}_X}(\mathcal{F}, \mathcal{G})$ sont eux aussi
quasi-cohérents, voir
Catégories dérivées des schémas, lemme
\ref{perfect-lemma-quasi-coherence-tensor-product}.
Le résultat suivant importe davantage pour ce qui nous occupe.

\begin{lemma}
\label{lemma-tensor-coherent}
Soit $X$ un schéma localement noethérien.
\begin{enumerate}
\item Si $\mathcal{F}$ et $\mathcal{G}$ sont des $\mathcal{O}_X$-modules cohérents,
alors $\text{Tor}_p^{\mathcal{O}_X}(\mathcal{F}, \mathcal{G})$ l'est aussi.
\item Si $L$ et $K$ appartiennent à $D^-_{\textit{Coh}}(\mathcal{O}_X)$, alors
il en est de même de $L \otimes_{\mathcal{O}_X}^\mathbf{L} K$.
\end{enumerate}
\end{lemma}

\begin{proof}
Expliquons comment démontrer (1) de manière plus élémentaire, et la partie (2)
à l'aide de la théorie générale développée précédemment.

\medskip\noindent
Démonstration de (1). Comme la formation de $\text{Tor}$ commute à la localisation,
nous pouvons supposer $X$ affine. Ainsi $X = \Spec(A)$ pour un anneau noethérien
$A$, et $\mathcal{F}$, $\mathcal{G}$ correspondent à des $A$-modules de type fini
$M$ et $N$ (Cohomologie des schémas, lemme
\ref{coherent-lemma-coherent-Noetherian}).
D'après Catégories dérivées des schémas, lemme
\ref{perfect-lemma-quasi-coherence-tensor-product}, nous pouvons
calculer les $\text{Tor}$ en calculant d'abord les $\text{Tor}$
de $M$ et $N$ sur $A$, puis en prenant le $\mathcal{O}_X$-module associé.
Comme les modules $\text{Tor}_p^A(M, N)$ sont de type fini d'après
Algèbre, lemme \ref{algebra-lemma-tor-noetherian},
nous concluons.

\medskip\noindent
D'après Catégories dérivées des schémas, lemme
\ref{perfect-lemma-identify-pseudo-coherent-noetherian},
l'hypothèse équivaut à demander que $L$ et $K$ soient
(localement) pseudo-cohérents. Alors $L \otimes_{\mathcal{O}_X}^\mathbf{L} K$
est pseudo-cohérent d'après
Cohomologie, lemme \ref{cohomology-lemma-tensor-pseudo-coherent}.
\end{proof}

\begin{lemma}
\label{lemma-compute-tor-nonsingular}
{\emergencystretch=2em Soit $X$ une variété non singulière.
Soient $\mathcal{F}$, $\mathcal{G}$ des $\mathcal{O}_X$-modules cohérents.
Le $\mathcal{O}_X$-module
$\text{Tor}_p^{\mathcal{O}_X}(\mathcal{F}, \mathcal{G})$
est cohérent, a pour fibre en $x$
$\text{Tor}_p^{\mathcal{O}_{X, x}}(\mathcal{F}_x, \mathcal{G}_x)$,
a son support contenu dans
$\text{Supp}(\mathcal{F}) \cap \text{Supp}(\mathcal{G})$, et
n'est non nul que pour $p \in \{0, \ldots, \dim(X)\}$.\par}
\end{lemma}

\begin{proof}
Le résultat sur les fibres a été discuté ci-dessus et implique la condition
sur le support. Les $\text{Tor}$ sont cohérents d'après
le lemme \ref{lemma-tensor-coherent}. L'annulation des $\text{Tor}$
en degrés négatifs résulte immédiatement de la construction. L'annulation
de $\text{Tor}_p$ pour $p > \dim(X)$ se voit comme suit :
les anneaux locaux $\mathcal{O}_{X, x}$ sont réguliers
(puisque $X$ est non singulière) et de dimension $\leq \dim(X)$
(Algèbre, lemme \ref{algebra-lemma-dimension-prime-polynomial-ring});
ainsi $\mathcal{O}_{X, x}$ est de dimension globale finie $\leq \dim(X)$
(Algèbre, lemme \ref{algebra-lemma-finite-gl-dim-finite-dim-regular}),
ce qui implique que les groupes $\text{Tor}$ de modules s'annulent au-delà de la dimension
(Compléments d'algèbre, lemme \ref{more-algebra-lemma-finite-gl-dim-tor-dimension}).
\end{proof}

\noindent
Soit $X$ une variété non singulière, et soient $W, V \subset X$
des sous-variétés fermées telles que $\dim(W) = s$ et $\dim(V) = r$.
Supposons que $V$ et $W$ se coupent proprement.
Dans ce cas, le lemme \ref{lemma-intersect-in-smooth} nous dit que toutes les composantes
irréductibles de $V \cap W$ ont une dimension égale à $r + s - \dim(X)$.
Les faisceaux $\text{Tor}_j^{\mathcal{O}_X}(\mathcal{O}_W, \mathcal{O}_V)$ sont
cohérents, à support dans $V \cap W$, et nuls si $j < 0$ ou $j > \dim(X)$
(lemme \ref{lemma-compute-tor-nonsingular}).
Nous définissons le {\it produit d'intersection} par
$$
W \cdot V = \sum\nolimits_i (-1)^i
[\text{Tor}_i^{\mathcal{O}_X}(\mathcal{O}_W, \mathcal{O}_V)]_{r + s - \dim(X)}.
$$
Soulignons que ceci n'a de sens qu'en vertu de notre hypothèse que
$V$ et $W$ se coupent proprement. Ce fait rendra nécessaire un lemme de
déplacement afin de définir le produit d'intersection en général.

\medskip\noindent
Avec cette notation, le cycle $V \cdot W$ est une combinaison linéaire
formelle $\sum e_Z Z$ des composantes irréductibles $Z$
de l'intersection $V \cap W$. Les entiers $e_Z$ sont appelés
les {\it multiplicités d'intersection}
$$
e_Z = e(X, V \cdot W, Z) =
\sum\nolimits_i
(-1)^i
\text{length}_{\mathcal{O}_{X, Z}}
\text{Tor}_i^{\mathcal{O}_{X, Z}}(\mathcal{O}_{W, Z}, \mathcal{O}_{V, Z})
$$
où $\mathcal{O}_{X, Z}$, resp.\ $\mathcal{O}_{W, Z}$,
resp.\ $\mathcal{O}_{V, Z}$ désigne l'anneau local de $X$, resp.\ $W$,
resp.\ $V$ au point générique de $Z$.
Ces sommes alternées des longueurs des $\text{Tor}$ satisfont à de nombreuses bonnes
propriétés, comme nous le verrons plus loin.

\medskip\noindent
Dans le cas d'intersections transversales, le nombre d'intersection vaut $1$.

\begin{lemma}
\label{lemma-transversal}
Soit $X$ une variété non singulière. Soient $V, W \subset X$ des
sous-variétés fermées qui se coupent proprement. Soit $Z$ une composante irréductible
de $V \cap W$, et supposons que la multiplicité
(au sens de la section \ref{section-cycle-of-closed}) de $Z$
dans le sous-schéma fermé $V \cap W$ vaille $1$.
Alors $e(X, V \cdot W, Z) = 1$, et $V$ et $W$ sont lisses
en un point général de $Z$.
\end{lemma}

\begin{proof}
Soit $(A, \mathfrak m, \kappa) =
(\mathcal{O}_{X, \xi}, \mathfrak m_\xi, \kappa(\xi))$, où $\xi \in Z$
est le point générique. Alors $\dim(A) = \dim(X) - \dim(Z)$, voir
Variétés, lemme \ref{varieties-lemma-dimension-locally-algebraic}.
Soient $I, J \subset A$ les idéaux définissant les traces de $V$ et $W$
dans $\Spec(A)$. Posons $\overline{I} = I + \mathfrak m^2/\mathfrak m^2$.
Alors $\dim_\kappa \overline{I} \leq \dim(X) - \dim(V)$, avec égalité
si et seulement si $A/I$ est régulier (cela résulte du lemme cité
ci-dessus et de la définition des anneaux réguliers, voir
Algèbre, définition \ref{algebra-definition-regular-local}
et la discussion qui la précède). De même pour $\overline{J}$.
Si la multiplicité vaut $1$, alors
$\text{length}_A(A/I + J) = 1$, donc $I + J = \mathfrak m$, et donc
$\overline{I} + \overline{J} = \mathfrak m/\mathfrak m^2$.
On a alors égalité partout (puisque l'intersection est
propre). Nous trouvons donc $f_1, \ldots, f_a \in I$ et $g_1, \ldots g_b \in J$
tels que $\overline{f}_1, \ldots, \overline{g}_b$ soit une base
de $\mathfrak m/\mathfrak m^2$. Alors $f_1, \ldots, g_b$ est un
système régulier de paramètres et une suite régulière
(Algèbre, lemme \ref{algebra-lemma-regular-ring-CM}).
Le même lemme montre que $A/(f_1, \ldots, f_a)$ est un anneau local régulier
de dimension $\dim(X) - \dim(V)$; ainsi $A/(f_1, \ldots, f_a) \to A/I$
est un isomorphisme (si le noyau est non nul, alors la dimension
de $A/I$ est strictement plus petite, voir
Algèbre, lemmes \ref{algebra-lemma-regular-domain} et
\ref{algebra-lemma-one-equation}).
Nous concluons que $I = (f_1, \ldots, f_a)$ et $J = (g_1, \ldots, g_b)$
par symétrie. Ainsi, le complexe de Koszul $K_\bullet(A, f_1, \ldots, f_a)$
sur $f_1, \ldots, f_a$ est une résolution de $A/I$, voir
Compléments d'algèbre, lemme \ref{more-algebra-lemma-regular-koszul-regular}.
Par conséquent,
\begin{align*}
\text{Tor}_p^A(A/I, A/J)
& =
H_p(K_\bullet(A, f_1, \ldots, f_a) \otimes_A A/J) \\
& =
H_p(K_\bullet(A/J, f_1 \bmod J, \ldots, f_a \bmod J))
\end{align*}
Puisque nous avons vu ci-dessus que $f_1 \bmod J, \ldots, f_a \bmod J$ est
un système régulier de paramètres de l'anneau local régulier $A/J$,
nous concluons qu'il n'y a qu'un seul groupe de cohomologie, à savoir
$H_0 = A/(I + J) = \kappa$. Ceci achève la démonstration.
\end{proof}

\begin{example}
\label{example-naive-multiplicity-wrong}
Dans cet exemple, nous montrons qu'il est nécessaire d'utiliser les
Tor supérieurs dans la formule des multiplicités d'intersection ci-dessus.
Soit $X$ une variété non singulière de dimension $4$.
Soit $p \in X$ un point fermé. Soient $V, W \subset X$
des sous-variétés fermées de $X$. Supposons qu'il existe un
isomorphisme
$$
\mathcal{O}_{X, p}^\wedge \cong \mathbf{C}[[x, y, z, w]]
$$
tel que l'idéal de $V$ soit $(xz, xw, yz, yw)$ et l'idéal
de $W$ soit $(x - z, y - w)$. Un calcul montre alors que
$$
\text{length}\ \mathbf{C}[[x, y, z, w]]/
(xz, xw, yz, yw, x - z, y - w) = 3
$$
D'autre part, la multiplicité $e(X, V \cdot W, p) = 2$,
comme on le voit du fait qu'au niveau formel local, $V$ est la
réunion de deux plans lisses $x = y = 0$ et $z = w = 0$ en $p$,
dont chacun a une multiplicité d'intersection $1$ avec le plan
$x - z = y - w = 0$ (lemme \ref{lemma-transversal}). Pour obtenir un
véritable exemple, prenons
un morphisme général $f : \mathbf{P}^2 \to \mathbf{P}^4$ donné par
$5$ polynômes homogènes de degré $> 1$. L'image
$V \subset \mathbf{P}^4 = X$ aura des singularités du type
décrit ci-dessus, car il existera $p_1, p_2 \in \mathbf{P}^2$
tels que $f(p_1) = f(p_2)$. Pour trouver $W$, prenons un plan général passant
par un tel point.
\end{example}






\section{Multiplicités algébriques}
\label{section-multiplicities}

\noindent
Soit $(A, \mathfrak m, \kappa)$ un anneau local noethérien.
Soit $M$ un $A$-module de type fini, et soit $I \subset A$ un idéal
de définition (Algèbre, définition \ref{algebra-definition-ideal-definition}).
Rappelons que la fonction
$$
\chi_{I, M}(n) = \text{length}_A(M/I^nM) =
\sum\nolimits_{p = 0, \ldots, n - 1} \text{length}_A(I^pM/I^{p + 1}M)
$$
est un polynôme numérique
(Algèbre, proposition \ref{algebra-proposition-hilbert-function-polynomial}).
Le degré de ce polynôme est égal à $\dim(\text{Supp}(M))$ d'après
Algèbre, lemme \ref{algebra-lemma-support-dimension-d}.

\begin{definition}
\label{definition-multiplicity}
Dans la situation ci-dessus, supposons $d \geq \dim(\text{Supp}(M))$.
Dans ce cas, si $d > \dim(\text{Supp}(M))$, nous posons $e_I(M, d) = 0$,
et si $d = \dim(\text{Supp}(M))$, nous posons $e_I(M, d)$ égal à $d!$
fois le coefficient dominant du polynôme numérique $\chi_{I, M}$.
Ainsi, dans les deux cas, nous avons
$$
\chi_{I, M}(n) \sim e_I(M, d) \frac{n^d}{d!} + \text{lower order terms}
$$
La {\it multiplicité de $M$ pour l'idéal de définition $I$}
est $e_I(M) = e_I(M, \dim(\text{Supp}(M)))$.
\end{definition}

\noindent
Nous avons les propriétés suivantes pour ces multiplicités.

\begin{lemma}
\label{lemma-multiplicity-ses}
Soit $A$ un anneau local noethérien. Soit $I \subset A$ un idéal de
définition. Soit $0 \to M' \to M \to M'' \to 0$ une suite exacte courte
de $A$-modules de type fini. Soit $d \geq \dim(\text{Supp}(M))$. Alors
$$
e_I(M, d) = e_I(M', d) + e_I(M'', d)
$$
\end{lemma}

\begin{proof}
Cela résulte immédiatement des définitions et de
Algèbre, lemme \ref{algebra-lemma-hilbert-ses-chi}.
\end{proof}

\begin{lemma}
\label{lemma-multiplicity-as-a-sum}
Soit $A$ un anneau local noethérien. Soit $I \subset A$ un idéal de
définition. Soit $M$ un $A$-module de type fini. Soit $d \geq \dim(\text{Supp}(M))$.
Alors
$$
e_I(M, d) =
\sum \text{length}_{A_\mathfrak p}(M_\mathfrak p) e_I(A/\mathfrak p, d)
$$
où la somme porte sur les idéaux premiers $\mathfrak p \subset A$ tels que
$\dim(A/\mathfrak p) = d$.
\end{lemma}

\begin{proof}
Le membre de gauche et le membre de droite sont tous deux additifs dans les suites
exactes courtes de modules de dimension $\leq d$, voir
le lemme \ref{lemma-multiplicity-ses} et
Algèbre, lemme \ref{algebra-lemma-length-additive}.
Ainsi, d'après Algèbre, lemme \ref{algebra-lemma-filter-Noetherian-module},
il suffit de démontrer ceci lorsque $M = A/\mathfrak q$ pour un
idéal premier $\mathfrak q$ de $A$ tel que $\dim(A/\mathfrak q) \leq d$.
Ce cas est évident.
\end{proof}

\begin{lemma}
\label{lemma-leading-coefficient}
Soit $P$ un polynôme de degré $r$ et de coefficient dominant $a$.
Alors
$$
r! a = \sum\nolimits_{i = 0, \ldots, r} (-1)^i{r \choose i} P(t - i)
$$
pour tout $t$.
\end{lemma}

\begin{proof}
Notons $\Delta$ l'opérateur qui, à un polynôme $P$, associe
le polynôme $\Delta(P) = P(t) - P(t - 1)$. Nous affirmons que
$$
\Delta^r(P) = \sum\nolimits_{i = 0, \ldots, r} (-1)^i {r \choose i} P(t - i)
$$
C'est vrai pour $r = 0, 1$ par inspection. Supposons que ce soit vrai pour $r$.
Alors nous calculons
\begin{align*}
\Delta^{r + 1}(P)
& =
\sum\nolimits_{i = 0, \ldots, r} (-1)^i {r \choose i} \Delta(P)(t - i) \\
& =
\sum\nolimits_{n = -r, \ldots, 0} (-1)^i {r \choose i}
(P(t - i) - P(t - i - 1))
\end{align*}
L'affirmation résulte donc de l'égalité
$$
{r + 1 \choose i} = {r \choose i} + {r \choose i - 1}
$$
Le lemme résulte du fait que $\Delta(P)$ est de degré $r - 1$
et de coefficient dominant $ra$ si le degré de $P$ est $r$.
\end{proof}

\noindent
Un fait important est que l'on peut calculer la multiplicité à l'aide
du complexe de Koszul. Rappelons que, si $R$ est un anneau et si
$f_1, \ldots, f_r \in R$, alors $K_\bullet(f_1, \ldots, f_r)$
désigne le complexe de Koszul, voir
Compléments d'algèbre, section \ref{more-algebra-section-koszul}.

\begin{theorem}
\label{theorem-multiplicity-with-koszul}
\begin{reference}
\cite[Theorem 1 in part B of Chapter IV]{Serre_algebre_locale}
\end{reference}
Soit $A$ un anneau local noethérien. Soit $I = (f_1, \ldots, f_r) \subset A$
un idéal de définition. Soit $M$ un $A$-module de type fini. Alors
$$
e_I(M, r) = \sum
(-1)^i\text{length}_A H_i(K_\bullet(f_1, \ldots, f_r) \otimes_A M)
$$
\end{theorem}

\begin{proof}
Transformons le complexe de Koszul $K_\bullet(f_1, \ldots, f_r)$ en un complexe de
cochaînes $K^\bullet$ en posant $K^n = K_{-n}(f_1, \ldots, f_r)$.
Alors $K^\bullet$ est placé en degrés $-r, \ldots, 0$ et
$H^i(K^\bullet \otimes_A M) = H_{-i}(K_\bullet(f_1, \ldots, f_r) \otimes_A M)$.
L'énoncé du théorème a un sens, car les modules
$H^i(K^\bullet \otimes M)$ sont annulés par $f_1, \ldots, f_r$
(Compléments d'algèbre, lemme \ref{more-algebra-lemma-homotopy-koszul})
et sont donc de longueur finie.
Définissons une filtration du complexe $K^\bullet$ en posant
$$
F^p(K^n \otimes_A M) =
I^{\max(0, p + n)}(K^n \otimes_A M),\quad p \in \mathbf{Z}
$$
Comme $f_i I^p \subset I^{p + 1}$, c'est une filtration par sous-complexes.
Nous avons ainsi un complexe filtré et obtenons une suite spectrale, voir
Homologie, section \ref{homology-section-filtered-complex}.
Nous avons
$$
E_0 = \bigoplus\nolimits_{p, q} E_0^{p, q} =
\bigoplus\nolimits_{p, q} \text{gr}^p(K^{p + q} \otimes_A M) =
\text{Gr}_I(K^\bullet \otimes_A M)
$$
Comme $K^n$ est libre de type fini, nous avons
$$
\text{Gr}_I(K^\bullet \otimes_A M) =
\text{Gr}_I(K^\bullet) \otimes_{\text{Gr}_I(A)} \text{Gr}_I(M)
$$
Remarquons que $\text{Gr}_I(K^\bullet)$ est le complexe de Koszul
sur $\text{Gr}_I(A)$ associé aux éléments
$\overline{f}_1, \ldots, \overline{f}_r \in I/I^2$.
Un calcul simple (omis)
montre que la différentielle $d_0$ sur $E_0$
coïncide avec la différentielle provenant du complexe de Koszul.
Comme $\text{Gr}_I(M)$ est un $\text{Gr}_I(A)$-module de type fini
et comme $\text{Gr}_I(A)$ est noethérien (en tant que quotient
de $A/I[x_1, \ldots, x_r]$ par $x_i \mapsto \overline{f}_i$), le
module de cohomologie $E_1 = \bigoplus E_1^{p, q}$
est un $\text{Gr}_I(A)$-module de type fini. Toutefois, comme ci-dessus,
$E_1$ est annulé par $\overline{f}_1, \ldots, \overline{f}_r$.
Nous concluons que $E_1$ est de longueur finie.
En particulier, nous trouvons que $\text{Gr}^p_F(K^\bullet \otimes M)$ est
acyclique pour $p \gg 0$.

\medskip\noindent
Vérifions ensuite que la suite spectrale ci-dessus converge
en utilisant Homologie, lemme \ref{homology-lemma-filtered-complex-ss-converges}.
Les égalités requises résultent aisément du lemme d'Artin-Rees
sous la forme énoncée dans Algèbre, lemme \ref{algebra-lemma-map-AR}.
Nous voyons ainsi que
\begin{align*}
\sum (-1)^i\text{length}_A(H^i(K^\bullet \otimes_A M))
& =
\sum (-1)^{p + q} \text{length}_A(E_\infty^{p, q}) \\
& =
\sum (-1)^{p + q} \text{length}_A(E_1^{p, q})
\end{align*}
car, comme nous l'avons vu ci-dessus, la longueur de $E_1$ est finie
(cela utilise bien sûr l'additivité des longueurs). Choisissons $t$ assez
grand pour que $\text{Gr}^p_F(K^\bullet \otimes M)$
soit acyclique pour $p \geq t$ (voir ci-dessus). En utilisant
de nouveau l'additivité, nous voyons que
$$
\sum (-1)^{p + q} \text{length}_A(E_1^{p, q}) =
\sum\nolimits_n \sum\nolimits_{p \leq t}
(-1)^n \text{length}_A(\text{gr}^p(K^n \otimes_A M))
$$
Ceci est égal à
$$
\sum\nolimits_{n = -r, \ldots, 0} (-1)^n{r \choose |n|} \chi_{I, M}(t + n)
$$
par le choix de la filtration ci-dessus et la définition de $\chi_{I, M}$ dans
Algèbre, section \ref{algebra-section-Noetherian-local}.
Le lemme résulte du lemme \ref{lemma-leading-coefficient}
et de la définition de $e_I(M, r)$.
\end{proof}

\begin{remark}[Généralisation triviale]
\label{remark-trivial-generalization}
Soit $(A, \mathfrak m, \kappa)$ un anneau local noethérien.
Soit $M$ un $A$-module de type fini. Soit $I \subset A$ un idéal.
Les conditions suivantes sont équivalentes :
\begin{enumerate}
\item $I' = I + \text{Ann}(M)$ est un idéal de définition
(Algèbre, définition \ref{algebra-definition-ideal-definition}),
\item l'image $\overline{I}$ de $I$ dans $\overline{A} = A/\text{Ann}(M)$
est un idéal de définition,
\item $\text{Supp}(M/IM) \subset \{\mathfrak m\}$,
\item $\dim(\text{Supp}(M/IM)) \leq 0$, et
\item $\text{length}_A(M/IM) < \infty$.
\end{enumerate}
Ceci résulte d'Algèbre, lemme \ref{algebra-lemma-support-point}
(détails omis). Si ces conditions sont satisfaites, on a $M/I^nM = M/(I')^nM$
pour tout $n$, et $M/I^nM = M/\overline{I}^nM$ pour tout $n$
si l'on considère $M$ comme un $\overline{A}$-module.
Nous pouvons donc définir
$$
\chi_{I, M}(n) = \text{length}_A(M/I^nM) =
\sum\nolimits_{p = 0, \ldots, n - 1} \text{length}_A(I^pM/I^{p + 1}M)
$$
et nous obtenons
$$
\chi_{I, M}(n) = \chi_{I', M}(n) = \chi_{\overline{I}, M}(n)
$$
pour tout $n$ par les égalités ci-dessus.
Tous les résultats d'Algèbre, section \ref{algebra-section-Noetherian-local},
ainsi que tous les résultats de la présente section, ont des analogues dans ce cadre.
En particulier, nous pouvons définir les multiplicités $e_I(M, d)$ pour
$d \geq \dim(\text{Supp}(M))$, et nous avons
$$
\chi_{I, M}(n) \sim e_I(M, d) \frac{n^d}{d!} + \text{lower order terms}
$$
comme dans le cas où $I$ est un idéal de définition.
\end{remark}



\section{Calcul des multiplicités d'intersection}
\label{section-computing-intersection-multiplicities}

\noindent
Dans cette section, nous étudions certains cas où les multiplicités d'intersection
peuvent se calculer par d'autres moyens. Voici un premier exemple.

\begin{lemma}
\label{lemma-intersection-multiplicity-CM}
Soit $X$ une variété non singulière, et soient $W, V \subset X$ des
sous-variétés fermées qui se coupent proprement. Soit $Z$ une composante irréductible
de $V \cap W$, de point générique $\xi$. Supposons que $\mathcal{O}_{W, \xi}$
et $\mathcal{O}_{V, \xi}$ soient de Cohen-Macaulay. Alors
$$
e(X, V \cdot W, Z) =
\text{length}_{\mathcal{O}_{X, \xi}}(\mathcal{O}_{V \cap W, \xi})
$$
où $V \cap W$ est l'intersection schématique.
En particulier, si $V$ et $W$ sont toutes deux de Cohen-Macaulay, alors
$V \cdot W = [V \cap W]_{\dim(V) + \dim(W) - \dim(X)}$.
\end{lemma}

\begin{proof}
Posons $A = \mathcal{O}_{X, \xi}$, $B = \mathcal{O}_{V, \xi}$, et
$C = \mathcal{O}_{W, \xi}$. D'après Auslander-Buchsbaum
(Algèbre, proposition \ref{algebra-proposition-Auslander-Buchsbaum}),
nous pouvons trouver une résolution finie $F_\bullet \to B$ par des modules libres
de type fini, de longueur
$$
\text{depth}(A) - \text{depth}(B) =
\dim(A) - \dim(B) = \dim(C)
$$
La première égalité vient de ce que $A$ et $B$ sont de Cohen-Macaulay, et la seconde
de ce que $V$ et $W$ se coupent proprement. Alors $F_\bullet \otimes_A C$ est un
complexe de modules libres de type fini représentant $B \otimes_A^\mathbf{L} C$;
ses modules de cohomologie ont donc leur support dans $\{\mathfrak m_A\}$.
D'après le lemme d'acyclicité (Algèbre, lemme \ref{algebra-lemma-acyclic}),
qui s'applique puisque $C$ est de Cohen-Macaulay,
nous concluons que $F_\bullet \otimes_A C$ n'a de
cohomologie non nulle qu'en degré $0$. Ceci achève la démonstration.
\end{proof}

\begin{lemma}
\label{lemma-one-ideal-ci}
Soit $A$ un anneau local noethérien. Soit $I = (f_1, \ldots, f_r)$ un idéal
engendré par une suite régulière. Soit $M$ un $A$-module de type fini. Supposons que
$\dim(\text{Supp}(M/IM)) = 0$. Alors
$$
e_I(M, r) = \sum (-1)^i\text{length}_A(\text{Tor}_i^A(A/I, M))
$$
Ici, $e_I(M, r)$ est défini comme dans la remarque \ref{remark-trivial-generalization}.
\end{lemma}

\begin{proof}
Puisque $f_1, \ldots, f_r$ est une suite régulière, le complexe de Koszul
$K_\bullet(f_1, \ldots, f_r)$ est une résolution de $A/I$ sur $A$, voir
Compléments d'algèbre, lemme
\ref{more-algebra-lemma-noetherian-finite-all-equivalent}.
Ainsi, le membre de droite est égal à
$$
\sum (-1)^i\text{length}_A H_i(K_\bullet(f_1, \ldots, f_r) \otimes_A M)
$$
Le résultat découle alors immédiatement du
théorème \ref{theorem-multiplicity-with-koszul} si $I$ est un idéal
de définition. En général, nous remplaçons $A$ par $\overline{A} = A/\text{Ann}(M)$
et $f_1, \ldots, f_r$ par $\overline{f}_1, \ldots, \overline{f}_r$,
ce qui est permis parce que
$$
K_\bullet(f_1, \ldots, f_r) \otimes_A M =
K_\bullet(\overline{f}_1, \ldots, \overline{f}_r) \otimes_{\overline{A}} M
$$
Puisque $e_I(M, r) = e_{\overline{I}}(M, r)$, où
$\overline{I} = (\overline{f}_1, \ldots, \overline{f}_r) \subset \overline{A}$
est un idéal de définition, le résultat découle aussi dans ce cas du
théorème \ref{theorem-multiplicity-with-koszul}.
\end{proof}

\begin{lemma}
\label{lemma-multiplicity-with-lci}
Soit $X$ une variété non singulière. Soient $W,V \subset X$ des
sous-variétés fermées qui se coupent proprement. Soit $Z$ une composante irréductible
de $V \cap W$, de point générique $\xi$.
Supposons que l'idéal de $V$ dans $\mathcal{O}_{X, \xi}$ soit défini par
une suite régulière $f_1, \ldots, f_c \in \mathcal{O}_{X, \xi}$.
Alors $e(X, V\cdot W, Z)$ est égal à $c!$ fois le coefficient dominant du
polynôme de Hilbert
$$
t \mapsto \text{length}_{\mathcal{O}_{X, \xi}}
\mathcal{O}_{W, \xi}/(f_1, \ldots, f_c)^t,\quad t \gg 0.
$$
En particulier, ce coefficient est $> 0$.
\end{lemma}

\begin{proof}
L'égalité
$$
e(X, V\cdot W, Z) = e_{(f_1, \ldots, f_c)}(\mathcal{O}_{W, \xi}, c)
$$
résulte du lemme plus général \ref{lemma-one-ideal-ci}.
Pour voir que $e_{(f_1, \ldots, f_c)}(\mathcal{O}_{W, \xi}, c)$ est
$> 0$, ou, de manière équivalente, que $e_{(f_1, \ldots, f_c)}(\mathcal{O}_{W, \xi}, c)$
est le coefficient dominant du polynôme de Hilbert,
il suffit de montrer que la
dimension de $\mathcal{O}_{W, \xi}$ vaut $c$, car le degré du
polynôme de Hilbert est égal à la dimension d'après
Algèbre, proposition \ref{algebra-proposition-dimension}.
Écrivons $\dim(V) = r$, $\dim(W) = s$ et $\dim(X) = n$. Alors
$\dim(Z) = r + s - n$ puisque l'intersection est propre. Ainsi,
le degré de transcendance de $\kappa(\xi)$ sur $\mathbf{C}$ est
$r + s - n$, voir Algèbre, lemme
\ref{algebra-lemma-dimension-prime-polynomial-ring}.
Nous avons $r + c = n$ parce que $V$ est défini par une suite régulière
dans un voisinage de $\xi$, voir
Diviseurs, lemme \ref{divisors-lemma-Noetherian-scheme-regular-ideal},
et alors le lemme \ref{lemma-pullback-by-regular-immersion}
s'applique (par exemple). Ainsi,
$$
\dim(\mathcal{O}_{W, \xi}) = s - (r + s - n) = s - ((n - c) + s - n) = c
$$
la première égalité résultant d'Algèbre, lemme
\ref{algebra-lemma-dimension-at-a-point-finite-type-field}.
\end{proof}

\begin{lemma}
\label{lemma-multiplicity-with-effective-Cartier-divisor}
Dans le lemme \ref{lemma-multiplicity-with-lci}, supposons que $c = 1$, c'est-à-dire que $V$
est un diviseur de Cartier effectif. Alors
$$
e(X, V \cdot W, Z) =
\text{length}_{\mathcal{O}_{X, \xi}}
(\mathcal{O}_{W, \xi}/f_1\mathcal{O}_{W, \xi}).
$$
\end{lemma}

\begin{proof}
Dans ce cas, l'image de $f_1$ dans $\mathcal{O}_{W, \xi}$ est non nulle par
propreté de l'intersection, donc est un non-diviseur de zéro. De plus,
$\mathcal{O}_{W, \xi}$ est un anneau local noethérien intègre de dimension $1$.
Ainsi,
$$
\text{length}_{\mathcal{O}_{X, \xi}}
(\mathcal{O}_{W, \xi}/f_1^t\mathcal{O}_{W, \xi}) =
t \text{length}_{\mathcal{O}_{X, \xi}}
(\mathcal{O}_{W, \xi}/f_1\mathcal{O}_{W, \xi})
$$
pour tout $t \geq 1$, voir Algèbre, lemme \ref{algebra-lemma-ord-additive}.
Ceci démontre le lemme.
\end{proof}

\begin{lemma}
\label{lemma-multiplicity-lci-CM}
Dans le lemme \ref{lemma-multiplicity-with-lci}, supposons que
l'anneau local $\mathcal{O}_{W, \xi}$ soit de Cohen-Macaulay. Alors nous
avons
$$
e(X, V \cdot W, Z) =
\text{length}_{\mathcal{O}_{X, \xi}} (\mathcal{O}_{W, \xi}/
f_1\mathcal{O}_{W, \xi} + \ldots + f_c\mathcal{O}_{W, \xi}).
$$
\end{lemma}

\begin{proof}
Ceci résulte immédiatement du lemme \ref{lemma-intersection-multiplicity-CM}.
On peut aussi le déduire du lemme \ref{lemma-multiplicity-with-lci}.
En effet, d'après Algèbre, lemme \ref{algebra-lemma-reformulate-CM},
nous voyons que $f_1, \ldots, f_c$ est une suite régulière dans
$\mathcal{O}_{W, \xi}$. Alors
Algèbre, lemme \ref{algebra-lemma-regular-quasi-regular} montre que
$f_1, \ldots, f_c$ est une suite quasi-régulière.
Ceci implique facilement que la longueur de
$\mathcal{O}_{W, \xi}/(f_1, \ldots, f_c)^t$ vaut
$$
{c + t \choose c}
\text{longueur}_{\mathcal{O}_{X, \xi}} (\mathcal{O}_{W, \xi}/
f_1\mathcal{O}_{W, \xi} + \ldots + f_c\mathcal{O}_{W, \xi}).
$$
L'examen du coefficient dominant permet de conclure.
\end{proof}


\section{Produit d'intersection à l'aide de la formule des Tor}
\label{section-intersection-product}

\noindent
Soit $X$ une variété non singulière. Soit
$\alpha = \sum n_i [W_i]$ un $r$-cycle et
$\beta = \sum_j m_j [V_j]$ un $s$-cycle sur $X$.
Supposons que $\alpha$ et $\beta$ se coupent proprement, voir
la définition \ref{definition-proper-intersection}.
Dans ce cas, nous posons
$$
\alpha \cdot \beta = \sum\nolimits_{i,j} n_i m_j W_i \cdot V_j.
$$
où $W_i \cdot V_j$ est défini comme à la section \ref{section-tor-formula}.
Si $\beta = [V]$, où $V$ est une sous-variété fermée de dimension $s$,
nous écrivons parfois $\alpha \cdot \beta = \alpha \cdot V$.

\begin{lemma}
\label{lemma-rational-equivalence-and-intersection}
Soit $X$ une variété non singulière. Soient $a, b \in \mathbf{P}^1$
deux points fermés distincts. Soit $k \geq 0$.
\begin{enumerate}
\item Si $W \subset X \times \mathbf{P}^1$ est une sous-variété fermée
de dimension $k + 1$ qui coupe $X \times a$ proprement, alors
\begin{enumerate}
\item $[W_a]_k = W \cdot X \times a$ comme cycles sur $X \times \mathbf{P}^1$, et
\item $[W_a]_k = \text{pr}_{X, *}(W \cdot X \times a)$ comme cycles sur $X$.
\end{enumerate}
\item Soit $\alpha$ un $(k + 1)$-cycle sur $X \times \mathbf{P}^1$
qui coupe $X \times a$ et $X \times b$ proprement. Alors
$pr_{X,*}( \alpha \cdot X \times a - \alpha \cdot X \times b)$
est rationnellement équivalent à zéro.
\item {\emergencystretch=1em Réciproquement, tout $k$-cycle rationnellement
équivalent à $0$ est de cette forme.\par}
\end{enumerate}
\end{lemma}

\begin{proof}
Observons d'abord que $X \times a$ est un diviseur de Cartier effectif de
$X \times \mathbf{P}^1$ et que $W_a$ est l'intersection schématique
de $W$ avec $X \times a$. Par conséquent, l'égalité de (1)(a) résulte
immédiatement des définitions et du calcul de la multiplicité d'intersection
dans le cas d'un diviseur de Cartier donné dans le
lemme \ref{lemma-multiplicity-with-effective-Cartier-divisor}.
L'assertion (1)(b) est vraie parce que $W_a \to X \times \mathbf{P}^1 \to X$ envoie
isomorphiquement sur son image, ce qui est la manière dont nous avons considéré $W_a$
comme sous-schéma fermé de $X$ à la section \ref{section-rational-equivalence}.
Les assertions (2) et (3) sont des conséquences formelles de (1) et des définitions.
\end{proof}

\noindent
Pour des intersections transverses de sous-schémas fermés, la multiplicité
d'intersection vaut $1$.

\begin{lemma}
\label{lemma-transversal-subschemes}
Soit $X$ une variété non singulière. Soient $r, s \geq 0$ et
$Y, Z \subset X$ des sous-schémas fermés tels que $\dim(Y) \leq r$ et
$\dim(Z) \leq s$. Supposons que $[Y]_r = \sum n_i[Y_i]$ et
$[Z]_s = \sum m_j[Z_j]$ se coupent proprement.
Soit $T$ une composante irréductible de $Y_{i_0} \cap Z_{j_0}$
pour certains $i_0$ et $j_0$, et supposons que la multiplicité
(au sens de la section \ref{section-cycle-of-closed}) de $T$
dans le sous-schéma fermé $Y \cap Z$ soit $1$.
Alors
\begin{enumerate}
\item le coefficient de $T$ dans $[Y]_r \cdot [Z]_s$ est $1$,
\item $Y$ et $Z$ sont non singuliers au point générique de $T$,
\item $n_{i_0} = 1$, $m_{j_0} = 1$, et
\item $T$ n'est contenu dans aucun $Y_i$ ni aucun $Z_j$ pour $i \not = i_0$ et
$j \not = j_0$.
\end{enumerate}
\end{lemma}

\begin{proof}
Posons $n = \dim(X)$, $a = n - r$, $b = n - s$. Remarquons que
$\dim(T) = r + s - n = n - a - b$ puisque, par hypothèse, les
intersections sont propres. Posons $(A, \mathfrak m, \kappa) =
(\mathcal{O}_{X, \xi}, \mathfrak m_\xi, \kappa(\xi))$, où $\xi \in T$
est le point générique. Alors $\dim(A) = a + b$, voir
Variétés, lemme \ref{varieties-lemma-dimension-locally-algebraic}.
Soient $I_0, I, J_0, J \subset A$ les idéaux qui définissent les traces de
$Y_{i_0}$, $Y$, $Z_{j_0}$, $Z$ dans $\Spec(A)$.
Alors $\dim(A/I) = \dim(A/I_0) = b$ et $\dim(A/J) = \dim(A/J_0) = a$
d'après la même référence. Posons $\overline{I} = I + \mathfrak m^2/\mathfrak m^2$.
Alors $I \subset I_0 \subset \mathfrak m$ et
$J \subset J_0 \subset \mathfrak m$, et $I + J = \mathfrak m$.
Le lemme \ref{lemma-transversal} et sa démonstration montrent que
$I_0 = (f_1, \ldots, f_a)$ et $J_0 = (g_1, \ldots, g_b)$
où $f_1, \ldots, g_b$ est un système régulier de paramètres
de l'anneau local régulier $A$. Comme $I + J = \mathfrak m$, l'application
$$
I \oplus J \to
\mathfrak m/\mathfrak m^2 =
\kappa f_1
\oplus \ldots \oplus
\kappa f_a
\oplus
\kappa g_1
\oplus \ldots \oplus
\kappa g_b
$$
est surjective. Nous en déduisons que l'on peut trouver
$f_1', \ldots, f_a' \in I$ et $g'_1, \ldots, g_b' \in J$
dont les classes résiduelles dans $\mathfrak m/\mathfrak m^2$ sont égales aux
classes résiduelles de $f_1, \ldots, f_a$ et de $g_1, \ldots, g_b$.
Alors $f'_1, \ldots, g'_b$ est un système régulier de paramètres de $A$.
Le lemme d'algèbre \ref{algebra-lemma-regular-ring-CM} montre que
$A/(f'_1, \ldots, f'_a)$ est un anneau local régulier de dimension $b$.
Ainsi, tout quotient non trivial de $A/(f'_1, \ldots, f'_a)$
est de dimension strictement inférieure
(Algèbre, lemmes \ref{algebra-lemma-regular-domain} et
\ref{algebra-lemma-one-equation}). Par conséquent, $I = (f'_1, \ldots, f'_a) = I_0$.
Par symétrie, $J = J_0$. Cela démontre (2), (3) et (4).
Enfin, le coefficient de $T$ dans $[Y]_r \cdot [Z]_s$
est le coefficient de $T$ dans $Y_{i_0} \cdot Z_{j_0}$, qui vaut
$1$ d'après le lemme \ref{lemma-transversal}.
\end{proof}



\section{Produit extérieur}
\label{section-exterior-product}

\noindent
Soient $X$ et $Y$ des variétés.
Soit $V$, resp.\ $W$, une sous-variété fermée de $X$, resp.\ de $Y$.
Le produit $V\times W$ est une sous-variété fermée de $X\times Y$
(lemme \ref{lemma-dimension-product-varieties}).
Pour un $k$-cycle $\alpha = \sum n_i [V_i]$ et un $l$-cycle
$\beta = \sum m_j [V_j]$ sur $Y$, nous définissons le
{\it produit extérieur} de $\alpha$ et $\beta$ comme le cycle
$\alpha \times \beta = \sum n_i m_j [V_i \times W_j]$.
Le produit extérieur définit une application $\mathbf{Z}$-linéaire
$$
Z_r(X) \otimes_\mathbf{Z} Z_s(Y) \longrightarrow Z_{r + s}(X \times Y)
$$
Montrons que le produit extérieur passe au quotient par l'équivalence rationnelle.

\begin{lemma}
\label{lemma-exterior-product-rational-equivalence}
Soient $X$ et $Y$ des variétés.
Soient $\alpha \in Z_r(X)$ et $\beta \in Z_s(Y)$.
Si $\alpha \sim_{rat} 0$ ou $\beta \sim_{rat} 0$, alors
$\alpha \times \beta \sim_{rat} 0$.
\end{lemma}

\begin{proof}
Par linéarité et symétrie en $X$ et $Y$, il suffit de le démontrer lorsque
$\alpha = [V]$ pour une sous-variété $V \subset X$ de dimension $r$, et
$\beta = [W_a]_s - [W_b]_s$ pour une sous-variété fermée
$W \subset Y \times \mathbf{P}^1$ de dimension $s + 1$ qui
coupe $Y \times a$ et $Y \times b$ proprement. Dans ce cas,
le lemme résulte de l'égalité
$$
[(V \times W)_a]_{r + s} = [V] \times [W_a]_s
$$
et de même en remplaçant $a$ par $b$. En effet, nous voyons alors que
$\alpha \times \beta = [(V \times W)_a]_{r + s} - [(V \times W)_b]_{r + s}$
comme voulu. Pour établir l'égalité affichée, remarquons l'égalité
$$
V \times W_a = (V \times W)_a
$$
de schémas. La projection $V \times W_a \to W_a$ induit une bijection
entre les composantes irréductibles (voir par exemple
Variétés, lemme \ref{varieties-lemma-bijection-irreducible-components}).
Soit $W' \subset W_a$ une composante irréductible de point générique $\zeta$.
Alors $V \times W'$ est la composante irréductible correspondante de
$V \times W_a$ (voir le lemme \ref{lemma-dimension-product-varieties}).
Soit $\xi$ le point générique de $V \times W'$. Nous devons montrer que
$$
\text{longueur}_{\mathcal{O}_{Y, \zeta}}(\mathcal{O}_{W_a, \zeta}) =
\text{longueur}_{\mathcal{O}_{X \times Y, \xi}}(
\mathcal{O}_{V \times W_a, \xi})
$$
Dans cette formule, nous pouvons remplacer
$\mathcal{O}_{Y, \zeta}$ par $\mathcal{O}_{W_a, \zeta}$ et
nous pouvons remplacer
$\mathcal{O}_{X \times Y, \zeta}$ par $\mathcal{O}_{V \times W_a, \zeta}$
(voir Algèbre, lemme \ref{algebra-lemma-length-independent}).
Puisque $\mathcal{O}_{W_a, \zeta} \to \mathcal{O}_{V \times W_a, \xi}$ est plat,
le lemme d'algèbre \ref{algebra-lemma-pullback-module} montre qu'il suffit
de vérifier que
$$
\text{longueur}_{\mathcal{O}_{V \times W_a, \xi}}(
\mathcal{O}_{V \times W_a, \xi}/
\mathfrak m_\zeta\mathcal{O}_{V \times W_a, \xi}) = 1
$$
C'est vrai, car le quotient de droite est l'anneau local
$\mathcal{O}_{V \times W', \xi}$ d'une variété en un point générique,
et est donc égal à $\kappa(\xi)$.
\end{proof}

\noindent
Nous en concluons que le produit extérieur définit un diagramme commutatif
$$
\xymatrix{
Z_r(X) \otimes_\mathbf{Z} Z_s(Y) \ar[r] \ar[d] &
Z_{r + s}(X \times Y) \ar[d] \\
\CH_r(X) \otimes_\mathbf{Z} \CH_s(Y) \ar[r] &
\CH_{r + s}(X \times Y)
}
$$
pour toute paire de variétés $X$ et $Y$. Pour les variétés non singulières,
nous pouvons considérer le produit extérieur comme un produit d'intersection
d'images réciproques.

\begin{lemma}
\label{lemma-exterior-product}
Soient $X$ et $Y$ des variétés non singulières.
Soient $\alpha \in Z_r(X)$ et $\beta \in Z_s(Y)$.
Alors
\begin{enumerate}
\item $\text{pr}_Y^*(\beta) = [X] \times \beta$ et
$\text{pr}_X^*(\alpha) = \alpha \times [Y]$,
\item $\alpha \times [Y]$ et $[X]\times \beta$
se coupent proprement sur $X\times Y$, et
\item on a
$\alpha \times \beta =
(\alpha \times [Y])\cdot ([X]\times\beta) =
pr_Y^*(\alpha) \cdot pr_X^*(\beta)$
dans $Z_{r + s}(X \times Y)$.
\end{enumerate}
\end{lemma}

\begin{proof}
Par linéarité, nous pouvons supposer que $\alpha = [V]$ et $\beta = [W]$.
Alors (1) affirme que $\text{pr}_Y^{-1}(W) = X \times W$ et
$\text{pr}_X^{-1}(V) = V \times Y$. Cela est clair.
L'assertion (2) est vraie parce que $X \times W \cap V \times Y = V \times W$ et
$\dim(V \times W) = r + s$ d'après le lemme \ref{lemma-dimension-product-varieties}.

\medskip\noindent
Démonstration de (3).
Soit $\xi$ le point générique de $V \times W$.
Puisque la projection $X \times W \to W$ est lisse comme changement de base du morphisme
$X \to \Spec(\mathbf{C})$, nous voyons que $X \times W$ est non singulier
en tout point situé au-dessus du point générique de $W$, en particulier en $\xi$.
Il en va de même pour $V \times Y$. Ainsi, $\mathcal{O}_{X \times W, \xi}$
et $\mathcal{O}_{V \times Y, \xi}$ sont des anneaux locaux de Cohen-Macaulay,
et le lemme \ref{lemma-intersection-multiplicity-CM} s'applique.
Comme $V \times Y \cap X \times W = V \times W$ schématiquement,
cela achève la démonstration.
\end{proof}



\section{Réduction à la diagonale}
\label{section-reduction-diagonal}

\noindent
Soit $X$ une variété non singulière. Nous utiliserons $\Delta$
pour désigner soit le morphisme diagonal $\Delta : X \to X \times X$,
soit l'image $\Delta \subset X \times X$.
La réduction à la diagonale est l'énoncé selon lequel
les produits d'intersection sur $X$ se ramènent aux produits d'intersection
des produits extérieurs avec la diagonale dans $X \times X$.

\begin{lemma}
\label{lemma-tor-and-diagonal}
Soit $X$ une variété non singulière.
\begin{enumerate}
\item Si $\mathcal{F}$ et $\mathcal{G}$ sont des $\mathcal{O}_X$-modules cohérents,
alors il existe des isomorphismes canoniques
$$
\text{Tor}_i^{\mathcal{O}_{X \times X}}(\mathcal{O}_\Delta,
\text{pr}_1^*\mathcal{F} \otimes_{\mathcal{O}_{X \times X}}
\text{pr}_2^*\mathcal{G})
=
\Delta_*\text{Tor}_i^{\mathcal{O}_X}(\mathcal{F}, \mathcal{G})
$$
\item Si $K$ et $M$ appartiennent à $D_\QCoh(\mathcal{O}_X)$, alors
il existe un isomorphisme canonique
$$
L\Delta^* \left(
L\text{pr}_1^*K \otimes_{\mathcal{O}_{X \times X}}^\mathbf{L} L\text{pr}_2^*M
\right)
= K \otimes_{\mathcal{O}_X}^\mathbf{L} M
$$
dans $D_\QCoh(\mathcal{O}_X)$ et un isomorphisme canonique
$$
\mathcal{O}_\Delta \otimes_{\mathcal{O}_{X \times X}}^\mathbf{L}
L\text{pr}_1^*K \otimes_{\mathcal{O}_{X \times X}}^\mathbf{L} L\text{pr}_2^*M
= \Delta_*(K \otimes_{\mathcal{O}_X}^\mathbf{L} M)
$$
dans $D_\QCoh(X \times X)$.
\end{enumerate}
\end{lemma}

\begin{proof}
Expliquons comment démontrer (1) de manière plus élémentaire et (2) à l'aide
d'une théorie plus générale. Comme (2) implique (1), le lecteur peut omettre la démonstration
de (1).

\medskip\noindent
Démonstration de (1). Choisissons un ouvert affine $\Spec(A) \subset X$.
Alors $A$ est une $\mathbf{C}$-algèbre noethérienne et
$\mathcal{F}$, $\mathcal{G}$ correspondent aux $A$-modules finis
$M$ et $N$ (Cohomologie des schémas, lemme
\ref{coherent-lemma-coherent-Noetherian}).
D'après Catégories dérivées de schémas, lemme
\ref{perfect-lemma-quasi-coherence-tensor-product}, nous pouvons
calculer $\text{Tor}_i$ sur $\mathcal{O}_X$
en calculant d'abord les $\text{Tor}$
de $M$ et $N$ sur $A$, puis en prenant le $\mathcal{O}_X$-module associé.
Pour les $\text{Tor}_i$ sur $\mathcal{O}_{X \times X}$, nous calculons
les Tor de $A$ et $M \otimes_\mathbf{C} N$ sur $A \otimes_\mathbf{C} A$
puis prenons le $\mathcal{O}_{X \times X}$-module associé.
Ainsi, sur cet ouvert affine, nous devons démontrer que
$$
\text{Tor}_i^{A \otimes_\mathbf{C} A}(A, M \otimes_\mathbf{C} N) =
\text{Tor}_i^A(M, N)
$$
Pour le voir, choisissons des résolutions $F_\bullet \to M$ et $G_\bullet \to M$
par des $A$-modules libres de type fini
(Algèbre, lemme \ref{algebra-lemma-resolution-by-finite-free}).
Remarquons que $\text{Tot}(F_\bullet \otimes_\mathbf{C} G_\bullet)$
est une résolution de $M \otimes_\mathbf{C} N$ puisqu'il calcule les
groupes de Tor sur $\mathbf{C}$ ! Bien entendu, les termes de
$F_\bullet \otimes_\mathbf{C} G_\bullet$ sont des
$A \otimes_\mathbf{C} A$-modules libres de type fini. Ainsi, le membre de gauche
de l'égalité affichée est le module
$$
H_i(A \otimes_{A \otimes_\mathbf{C} A}
\text{Tot}(F_\bullet \otimes_\mathbf{C} G_\bullet))
$$
et le membre de droite est le module
$$
H_i(\text{Tot}(F_\bullet \otimes_A G_\bullet))
$$
Comme $A \otimes_{A \otimes_\mathbf{C} A} (F_p \otimes_\mathbf{C} G_q)
= F_p \otimes_A G_q$, nous voyons que ces modules sont égaux.
Cela définit un isomorphisme sur l'ouvert affine $\Spec(A) \times \Spec(A)$
(ce qui suffit pour l'application à l'égalité des nombres d'intersection).
Nous omettons la démonstration du recollement de ces isomorphismes.

\medskip\noindent
Démonstration de (2). La seconde assertion résulte de la première par la
formule de projection telle qu'énoncée dans
Catégories dérivées de schémas, lemme \ref{perfect-lemma-cohomology-base-change}.
Pour établir la première, représentons $K$ et $M$ par des complexes K-plats
$\mathcal{K}^\bullet$ et $\mathcal{M}^\bullet$.
Comme l'image réciproque et le produit tensoriel préservent les complexes K-plats
(Cohomologie, lemmes \ref{cohomology-lemma-tensor-product-K-flat} et
\ref{cohomology-lemma-pullback-K-flat})
nous voyons qu'il suffit de montrer
$$
\Delta^*\text{Tot}(
\text{pr}_1^*\mathcal{K}^\bullet
\otimes_{\mathcal{O}_{X \times X}} \text{pr}_2^*\mathcal{M}^\bullet)
=
\text{Tot}(
\mathcal{K}^\bullet \otimes_{\mathcal{O}_X} \mathcal{M}^\bullet)
$$
Il suffit donc de voir qu'il existe des isomorphismes canoniques
$$
\Delta^*(\text{pr}_1^*\mathcal{K}
\otimes_{\mathcal{O}_{X \times X}} \text{pr}_2^*\mathcal{M})
\longrightarrow
\mathcal{K} \otimes_{\mathcal{O}_X} \mathcal{M}
$$
lorsque $\mathcal{K}$ et $\mathcal{M}$ sont des $\mathcal{O}_X$-modules
(pas nécessairement quasi-cohérents ni plats).
Nous omettons les détails.
\end{proof}

\begin{lemma}
\label{lemma-reduction-diagonal}
Soit $X$ une variété non singulière. Soient $\alpha$, resp.\ $\beta$,
un $r$-cycle, resp.\ un $s$-cycle sur $X$. Supposons que $\alpha$ et $\beta$
se coupent proprement. Alors
\begin{enumerate}
\item $\alpha \times \beta$ et $[\Delta]$ se coupent proprement,
\item on a $\Delta_*(\alpha \cdot \beta) = [\Delta] \cdot \alpha\times\beta$
comme cycles sur $X \times X$,
\item si $X$ est propre, alors
$\text{pr}_{1, *}([\Delta] \cdot \alpha\times\beta) = \alpha\cdot\beta$,
où $pr_1 : X\times X \to X$ est la projection.
\end{enumerate}
\end{lemma}

\begin{proof}
Par linéarité, il suffit de le démontrer lorsque $\alpha = [V]$ et $\beta = [W]$
pour des sous-variétés fermées $V \subset X$ et $W \subset Y$ qui se coupent
proprement. Rappelons que $V \times W$ est une sous-variété fermée de dimension $r + s$.
Remarquons que, schématiquement, on a
$V \cap W = \Delta^{-1}(V \times W)$ ainsi que
$\Delta(V \cap W) = \Delta \cap V \times W$.
Cela démontre (1).

\medskip\noindent
Démonstration de (2). Soit $Z \subset V \cap W$ une composante irréductible
de point générique $\xi$. Nous devons montrer que le coefficient de
$Z$ dans $\alpha \cdot \beta$ est égal au coefficient de
$\Delta(Z)$ dans $[\Delta] \cdot \alpha \times \beta$. Le premier est donné
par l'entier
$$
\sum (-1)^i
\text{longueur}_{\mathcal{O}_{X, \xi}}
\text{Tor}_i^{\mathcal{O}_X}(\mathcal{O}_V, \mathcal{O}_W)_\xi
$$
et le second par l'entier
$$
\sum (-1)^i
\text{longueur}_{\mathcal{O}_{X \times Y, \Delta(\xi)}}
\text{Tor}_i^{\mathcal{O}_{X \times Y}}(
\mathcal{O}_\Delta, \mathcal{O}_{V \times W})_{\Delta(\xi)}
$$
Cependant, d'après le lemme \ref{lemma-tor-and-diagonal}, on a
$$
\text{Tor}_i^{\mathcal{O}_X}(\mathcal{O}_V, \mathcal{O}_W)_\xi \cong
\text{Tor}_i^{\mathcal{O}_{X \times Y}}(
\mathcal{O}_\Delta, \mathcal{O}_{V \times W})_{\Delta(\xi)}
$$
comme $\mathcal{O}_{X \times X, \Delta(\xi)}$-modules. On a donc égalité
des longueurs (précisément, d'après Algèbre, lemme \ref{algebra-lemma-length-independent}
pour être précis).

\medskip\noindent
L'assertion (2) implique (3), car
$\text{pr}_{1, *} \circ \Delta_* = \text{id}$ d'après
le lemme \ref{lemma-compose-pushforward}.
\end{proof}

\begin{proposition}
\label{proposition-positivity}
\begin{reference}
C'est l'un des résultats principaux de \cite{Serre_algebre_locale}.
\end{reference}
Soit $X$ une variété non singulière. Soient $V \subset X$ et
$W \subset Y$ des sous-variétés fermées qui se coupent proprement.
Soit $Z \subset V \cap W$ une composante irréductible.
Alors $e(X, V \cdot W, Z) > 0$.
\end{proposition}

\begin{proof}
D'après le lemme \ref{lemma-reduction-diagonal}, on a
$$
e(X, V \cdot W, Z) = e(X \times X, \Delta \cdot V \times W, \Delta(Z))
$$
Puisque $\Delta : X \to X \times X$ est une immersion régulière
(voir le lemme \ref{lemma-diagonal-regular-immersion}), nous voyons que
$e(X \times X, \Delta \cdot V \times W, \Delta(Z))$ est un entier strictement
positif d'après le lemme \ref{lemma-multiplicity-with-lci}.
\end{proof}

\noindent
Le lemme suivant est un lemme clé du développement de la théorie
présenté dans ce chapitre. Essentiellement, ce lemme affirme que
les nombres d'intersection vérifient une propriété d'additivité convenable.

\begin{lemma}
\label{lemma-tor-sheaf}
\begin{reference}
\cite[Chapter V]{Serre_algebre_locale}
\end{reference}
Soit $X$ une variété non singulière. Soient $\mathcal{F}$ et
$\mathcal{G}$ des faisceaux cohérents sur $X$ tels que
$\dim(\text{Supp}(\mathcal{F})) \leq r$,
$\dim(\text{Supp}(\mathcal{G})) \leq s$, et
$\dim(\text{Supp}(\mathcal{F}) \cap \text{Supp}(\mathcal{G}) )
\leq r + s - \dim X$. Dans ce cas, $[\mathcal{F}]_r$ et $[\mathcal{G}]_s$
se coupent proprement et
$$
[\mathcal{F}]_r \cdot [\mathcal{G}]_s =
\sum (-1)^p
[\text{Tor}_p^{\mathcal{O}_X}(\mathcal{F}, \mathcal{G})]_{r + s - \dim(X)}.
$$
\end{lemma}

\begin{proof}
L'assertion selon laquelle $[\mathcal{F}]_r$ et $[\mathcal{G}]_s$ se coupent proprement
est immédiate. Comme nous démontrons une égalité de cycles, nous pouvons travailler
localement sur $X$. (Remarquons que la formation du produit d'intersection
des cycles, celle des faisceaux de $\text{Tor}$ et
celle du cycle associé à un faisceau cohérent commutent chacune à la
restriction aux sous-schémas ouverts.) Nous pouvons donc supposer, et supposons, $X$ affine.

\medskip\noindent
Posons
$$
RHS(\mathcal{F}, \mathcal{G}) = [\mathcal{F}]_r \cdot [\mathcal{G}]_s
\quad\text{et}\quad
LHS(\mathcal{F}, \mathcal{G}) = \sum (-1)^p
[\text{Tor}_p^{\mathcal{O}_X}(\mathcal{F}, \mathcal{G})]_{r + s - \dim(X)}
$$
Considérons une suite exacte
$$
0 \to \mathcal{F}_1 \to \mathcal{F}_2 \to \mathcal{F}_3 \to 0
$$
de faisceaux cohérents sur $X$ telle que
$\text{Supp}(\mathcal{F}_i) \subset \text{Supp}(\mathcal{F})$ ;
alors $LHS(\mathcal{F}_i, \mathcal{G})$ et
$RHS(\mathcal{F}_i, \mathcal{G})$ sont définis pour $i = 1, 2, 3$,
et l'on a
$$
RHS(\mathcal{F}_2, \mathcal{G}) =
RHS(\mathcal{F}_1, \mathcal{G}) + RHS(\mathcal{F}_3, \mathcal{G})
$$
et de même pour LHS. En effet, la condition sur les supports garantit que
tout est défini, et la suite exacte ainsi que l'additivité des longueurs
donnent
$$
[\mathcal{F}_2]_r = [\mathcal{F}_1]_r  + [\mathcal{F}_3]_r
$$
(Homologie de Chow, lemme \ref{chow-lemma-additivity-sheaf-cycle}),
ce qui entraîne l'additivité pour RHS. La suite exacte longue des $\text{Tor}$
$$
\ldots \to \text{Tor}_1(\mathcal{F}_3, \mathcal{G}) \to
\text{Tor}_0(\mathcal{F}_1, \mathcal{G}) \to
\text{Tor}_0(\mathcal{F}_2, \mathcal{G}) \to
\text{Tor}_0(\mathcal{F}_3, \mathcal{G}) \to 0
$$
et l'additivité des longueurs, comme précédemment, entraînent l'additivité pour LHS.

\medskip\noindent
{\emergencystretch=1em D'après Algèbre, lemme \ref{algebra-lemma-filter-Noetherian-module},
et puisque $X$ est affine, nous pouvons trouver une filtration de $\mathcal{F}$
dont les quotients gradués sont des faisceaux structuraux de sous-variétés fermées de
$\text{Supp}(\mathcal{F})$. L'additivité démontrée au paragraphe précédent
montre qu'il suffit de démontrer $LHS = RHS$ avec
$\mathcal{F}$ remplacé par $\mathcal{O}_V$, où
$V \subset \text{Supp}(\mathcal{F})$.
Par symétrie, nous pouvons faire de même pour $\mathcal{G}$.
Nous sommes ainsi ramenés à démontrer que\par}
$$
LHS(\mathcal{O}_V, \mathcal{O}_W) = RHS(\mathcal{O}_V, \mathcal{O}_W)
$$
où $W \subset \text{Supp}(\mathcal{G})$ est une sous-variété fermée.
Si $\dim(V) = r$ et $\dim(W) = s$, alors cette égalité est la
{\bf définition} de $V \cdot W$. En revanche, si
$\dim(V) < r$ ou $\dim(W) < s$, c'est-à-dire si $[V]_r = 0$ ou $[W]_s = 0$,
alors nous devons démontrer que $RHS(\mathcal{O}_V, \mathcal{O}_W) = 0$
\footnote{Le lecteur peut voir que ce n'est pas trivial en
prenant $r = s = 1$, $X$ une surface non singulière et $V = W$
un point fermé $x$ de $X$. Dans ce cas, il y a $3$
$\text{Tor}$ non nuls, de longueurs $1, 2, 1$ en $x$.}.

\medskip\noindent
Soit $Z \subset V \cap W$ une composante irréductible de dimension
$r + s - \dim(X)$. C'est la dimension maximale d'une composante,
et il suffit de montrer que le coefficient de $Z$ dans $RHS$ est nul.
Soit $\xi \in Z$ le point générique. Posons $A = \mathcal{O}_{X, \xi}$,
$B = \mathcal{O}_{X \times X, \Delta(\xi)}$, et
$C = \mathcal{O}_{V \times W, \Delta(\xi)}$.
D'après le lemme \ref{lemma-tor-and-diagonal}, on a
$$
\text{coefficient de }Z\text{ dans }
RHS(\mathcal{O}_V, \mathcal{O}_W) = 
\sum (-1)^i
\text{longueur}_B \text{Tor}_i^B(A, C)
$$
Comme $\dim(V) < r$ ou $\dim(W) < s$, on a $\dim(V \times W) < r + s$,
ce qui entraîne $\dim(C) < \dim(X)$ (on omet un petit détail). De plus, le
noyau $I$ de $B \to A$ est engendré par une suite régulière de
longueur $\dim(X)$ (lemme \ref{lemma-diagonal-regular-immersion}).
Il y a donc annulation d'après le lemme \ref{lemma-one-ideal-ci}, car
la fonction de Hilbert de $C$ relativement à $I$ est de degré $\dim(C) < n$
d'après Algèbre, proposition \ref{algebra-proposition-dimension}.
\end{proof}

\begin{remark}
\label{remark-Serre-conjectures}
Soit $(A, \mathfrak m, \kappa)$ un anneau local régulier.
Soient $M$ et $N$ des $A$-modules finis non nuls tels que $M \otimes_A N$
soit supporté par $\{\mathfrak m\}$. Alors
$$
\chi(M, N) = \sum (-1)^i \text{longueur}_A \text{Tor}_i^A(M, N)
$$
est fini. Posons $r = \dim(\text{Supp}(M))$ et $s = \dim(\text{Supp}(N))$.
Dans \cite{Serre_algebre_locale}, on démontre que $r + s \leq \dim(A)$
et les conjectures suivantes y sont formulées :
\begin{enumerate}
\item si $r + s < \dim(A)$, alors $\chi(M, N) = 0$, et
\item si $r + s = \dim(A)$, alors $\chi(M, N) > 0$.
\end{enumerate}
Les arguments qui démontrent le lemme \ref{lemma-tor-sheaf} et
la proposition \ref{proposition-positivity} peuvent être exploités
(comme dans le texte de Serre) pour montrer que (1) et (2) sont
vraies si $A$ contient un corps. À l'heure actuelle, la conjecture (1) est démontrée
en général, et l'on sait que $\chi(M, N) \geq 0$ en général (Gabber).
À notre connaissance, la positivité reste un problème ouvert.
\end{remark}



\section{Associativité des intersections}
\label{section-associative}

\noindent
Il est clair que les intersections propres définies ci-dessus sont commutatives.
À l'aide du lemme clé \ref{lemma-tor-sheaf}, nous pouvons démontrer que les produits
d'intersection (propres) sont associatifs.

\begin{lemma}
\label{lemma-associative}
Soit $X$ une variété non singulière. Soient $U, V, W$ des sous-variétés
fermées. Supposons que $U, V, W$ se coupent proprement deux à deux
et que $\dim(U \cap V \cap W) \leq \dim(U) + \dim(V) + \dim(W) - 2\dim(X)$.
Alors
$$
U \cdot (V \cdot W) = (U \cdot V) \cdot W
$$
comme cycles sur $X$.
\end{lemma}

\begin{proof}
Nous allons utiliser le lemme \ref{lemma-tor-sheaf} sans plus le mentionner.
Cela implique que
\begin{align*}
V \cdot W
& =
\sum (-1)^i [\text{Tor}_i(\mathcal{O}_V, \mathcal{O}_W)]_{b + c - n} \\
U \cdot (V \cdot W)
& =
\sum (-1)^{i + j}
[
\text{Tor}_j(\mathcal{O}_U, \text{Tor}_i(\mathcal{O}_V, \mathcal{O}_W))
]_{a + b + c - 2n} \\
U \cdot V
& =
\sum (-1)^i [\text{Tor}_i(\mathcal{O}_U, \mathcal{O}_V)]_{a + b - n} \\
(U \cdot V) \cdot W
& =
\sum (-1)^{i + j}
[
\text{Tor}_j(\text{Tor}_i(\mathcal{O}_U, \mathcal{O}_V), \mathcal{O}_W))
]_{a + b + c - 2n}
\end{align*}
où $\dim(U) = a$, $\dim(V) = b$, $\dim(W) = c$, $\dim(X) = n$.
Les hypothèses du lemme garantissent que les faisceaux cohérents
des formules ci-dessus vérifient les majorations requises pour les dimensions
des supports, de sorte que ces formules ont un sens. Considérons maintenant l'objet
$$
K =
\mathcal{O}_U \otimes^\mathbf{L}_{\mathcal{O}_X} \mathcal{O}_V
\otimes^\mathbf{L}_{\mathcal{O}_X} \mathcal{O}_W
$$
de la catégorie dérivée $D_{\textit{Coh}}(\mathcal{O}_X)$.
Nous affirmons que les expressions obtenues ci-dessus pour
$U \cdot (V \cdot W)$ et $(U \cdot V) \cdot W$
sont égaux à
$$
\sum (-1)^k [H^k(K)]_{a + b + c - 2n}
$$
Cela démontrera le lemme. Par symétrie, il suffit de démontrer l'une
de ces égalités. Pour cela, représentons $\mathcal{O}_U$
et $\mathcal{O}_V \otimes_{\mathcal{O}_X}^\mathbf{L} \mathcal{O}_W$
par des complexes K-plats $M^\bullet$ et $L^\bullet$ et utilisons la
suite spectrale associée au complexe double
$M^\bullet \otimes_{\mathcal{O}_X} L^\bullet$ dans
Homologie, section \ref{homology-section-double-complex}.
Il s'agit d'une suite spectrale dont la page $E_2$ est
$$
E_2^{p, q} =
\text{Tor}_{-p}(\mathcal{O}_U, \text{Tor}_{-q}(\mathcal{O}_V, \mathcal{O}_W))
$$
et qui converge vers $H^{p + q}(K)$ (détails omis ; comparer à
Compléments d'algèbre, exemple \ref{more-algebra-example-tor}).
Puisque les longueurs sont additives dans les suites
exactes courtes, le résultat en découle.
\end{proof}




\section{Image réciproque plate et produits d'intersection}
\label{section-flat-pullback-and-intersection-products}

\noindent
Brève discussion de l'interaction entre les intersections et
l'image réciproque plate.

\begin{lemma}
\label{lemma-flat-pull-back-and-intersections-sheaves}
Soit $f : X \to Y$ un morphisme plat de variétés non singulières. Posons
$e = \dim(X) - \dim(Y)$. Soient $\mathcal{F}$ et $\mathcal{G}$ des faisceaux
cohérents sur $Y$ tels que $\dim(\text{Supp}(\mathcal{F})) \leq r$,
$\dim(\text{Supp}(\mathcal{G})) \leq s$ et
$\dim(\text{Supp}(\mathcal{F}) \cap \text{Supp}(\mathcal{G}) )
\leq r + s - \dim(Y)$. Alors les cycles
$[f^*\mathcal{F}]_{r + e}$ et $[f^*\mathcal{G}]_{s + e}$
se coupent proprement et
$$
f^*([\mathcal{F}]_r \cdot [\mathcal{G}]_s) =
[f^*\mathcal{F}]_{r + e} \cdot [f^*\mathcal{G}]_{s + e}
$$
\end{lemma}

\begin{proof}
Le fait que $[f^*\mathcal{F}]_{r + e}$ et $[f^*\mathcal{G}]_{s + e}$
se coupent proprement résulte immédiatement de l'hypothèse que $f$ est de
dimension relative $e$. D'après les
lemmes \ref{lemma-tor-sheaf} et \ref{lemma-pullback},
il suffit de montrer que
$$
f^*\text{Tor}_i^{\mathcal{O}_Y}(\mathcal{F}, \mathcal{G}) =
\text{Tor}_i^{\mathcal{O}_X}(f^*\mathcal{F}, f^*\mathcal{G})
$$
comme $\mathcal{O}_X$-modules. Cela résulte de
Cohomologie, lemme \ref{cohomology-lemma-pullback-tensor-product},
et du fait que $f^*$ est exact, donc $Lf^*\mathcal{F} = f^*\mathcal{F}$,
et de même pour $\mathcal{G}$.
\end{proof}

\begin{lemma}
\label{lemma-flat-pullback-and-intersections}
Soit $f : X \to Y$ un morphisme plat de variétés non singulières.
Soit $\alpha$ un $r$-cycle sur $Y$ et $\beta$ un $s$-cycle sur $Y$.
Supposons que $\alpha$ et $\beta$ se coupent proprement. Alors $f^*\alpha$
et $f^*\beta$ se coupent proprement et
$f^*( \alpha \cdot \beta ) = f^*\alpha \cdot f^*\beta$.
\end{lemma}

\begin{proof}
Par linéarité, nous pouvons supposer que $\alpha = [V]$ et $\beta = [W]$
pour des sous-variétés fermées $V, W \subset Y$ de dimensions $r, s$.
Disons que $f$ est de dimension relative $e$. Le lemme est alors un cas particulier du
lemme \ref{lemma-flat-pull-back-and-intersections-sheaves},
car $[V] = [\mathcal{O}_V]_r$, $[W] = [\mathcal{O}_W]_r$,
$f^*[V] = [f^{-1}(V)]_{r + e} = [f^*\mathcal{O}_V]_{r + e}$ et
$f^*[W] = [f^{-1}(W)]_{s + e} = [f^*\mathcal{O}_W]_{s + e}$.
\end{proof}


\section{Formule de projection pour les morphismes plats propres}
\label{section-projection-formula-flat}

\noindent
Brève discussion de la formule de projection pour les morphismes plats propres.

\begin{lemma}
\label{lemma-projection-formula-flat}
\begin{reference}
Voir \cite[Chapter V, C), Section 7, formula (10)]{Serre_algebre_locale}
pour une formule plus générale.
\end{reference}
Soit $f : X \to Y$ un morphisme plat propre de variétés non singulières.
Posons $e = \dim(X) - \dim(Y)$. Soit $\alpha$ un $r$-cycle sur $X$ et soit
$\beta$ un $s$-cycle sur $Y$. Supposons que $\alpha$ et $f^*(\beta)$ se coupent
proprement. Alors $f_*(\alpha)$ et $\beta$ se coupent proprement et
$$
f_*(\alpha) \cdot \beta = f_*( \alpha \cdot f^*\beta)
$$
\end{lemma}

\begin{proof}
Par linéarité, on se ramène au cas où $\alpha = [V]$ et
$\beta = [W]$ pour des sous-variétés fermées $V \subset X$ et
$W \subset Y$ de dimensions $r$ et $s$. Alors $f^{-1}(W)$ est de
dimension pure $s + e$. Nous supposons que les cycles
$[V]$ et $f^*[W]$ se coupent proprement. Nous utiliserons sans
autre mention le fait que $V \cap f^{-1}(W) \to f(V) \cap W$
est surjectif.

\medskip\noindent
Soit $a$ la dimension de la fibre générique de $V \to f(V)$.
Si $a > 0$, alors $f_*[V] = 0$. En particulier, $f_*\alpha$ et $\beta$
se coupent proprement. Pour achever ce cas, il faut montrer que
$f_*([V] \cdot f^*[W]) = 0$. Cependant, puisque toute fibre de
$V \to f(V)$ est de dimension $\geq a$ (voir
Morphismes, lemme \ref{morphisms-lemma-openness-bounded-dimension-fibres}),
nous concluons que toute composante irréductible $Z$ de $V \cap f^{-1}(W)$
a des fibres de dimension $\geq a$ au-dessus de $f(Z)$. Cela implique
assurément l'assertion voulue.

\medskip\noindent
Supposons que $V \to f(V)$ soit génériquement fini. Soit $Z \subset f(V) \cap W$
une composante irréductible. Soient $Z_i \subset V \cap f^{-1}(W)$,
$i = 1, \ldots, t$, les composantes irréductibles de $V \cap f^{-1}(W)$
qui dominent $Z$. Par hypothèse, chaque $Z_i$ est de dimension
$r + s + e - \dim(X) = r + s - \dim(Y)$. Par conséquent,
$\dim(Z) \leq r + s - \dim(Y)$. On voit donc que $f(V)$ et $W$
se coupent proprement, que $\dim(Z) = r + s - \dim(Y)$ et que chaque
$Z_i \to Z$ est génériquement fini. Il s'ensuit notamment que
$V \to f(V)$ a une fibre finie au point générique $\xi$ de $Z$.
Ainsi, $V \to Y$ est fini au-dessus d'un voisinage ouvert de $\xi$, voir
Cohomologie des schémas, lemme
\ref{coherent-lemma-proper-finite-fibre-finite-in-neighbourhood}.
En utilisant une formule de projection très générale pour les produits tensoriels dérivés, on obtient
$$
Rf_*(\mathcal{O}_V \otimes_{\mathcal{O}_X}^\mathbf{L} Lf^*\mathcal{O}_W) =
Rf_*\mathcal{O}_V \otimes_{\mathcal{O}_Y}^\mathbf{L} \mathcal{O}_W
$$
{\emergencystretch=1em voir Catégories dérivées de schémas, lemme
\ref{perfect-lemma-cohomology-base-change}.
Comme $f$ est plat, on a $Lf^*\mathcal{O}_W = f^*\mathcal{O}_W$.
Puisque $f|_V$ est fini au-dessus d'un voisinage ouvert de $\xi$, on a\par}
$$
(Rf_*\mathcal{F})_\xi = (f_*\mathcal{F})_\xi
$$
pour tout faisceau cohérent sur $X$ dont le support est contenu dans $V$
(voir Cohomologie des schémas, lemme
\ref{coherent-lemma-higher-direct-images-zero-finite-fibre}). Nous
en concluons que
\begin{equation}
\label{equation-stalks}
\left(
f_*\text{Tor}_i^{\mathcal{O}_X}(\mathcal{O}_V, f^*\mathcal{O}_W)
\right)_\xi =
\left(\text{Tor}_i^{\mathcal{O}_Y}(f_*\mathcal{O}_V, \mathcal{O}_W)\right)_\xi
\end{equation}
pour tout $i$. Puisque $f^*[W] = [f^*\mathcal{O}_W]_{s + e}$ d'après le
lemme \ref{lemma-pullback}, on a
$$
[V] \cdot f^*[W] =
\sum (-1)^i
[\text{Tor}_i^{\mathcal{O}_X}(\mathcal{O}_V,
f^*\mathcal{O}_W)]_{r + s - \dim(Y)}
$$
d'après le lemme \ref{lemma-tor-sheaf}. En appliquant le
lemme \ref{lemma-push-coherent},
on trouve
$$
f_*([V] \cdot f^*[W]) =
\sum (-1)^i
[f_*\text{Tor}_i^{\mathcal{O}_X}(\mathcal{O}_V,
f^*\mathcal{O}_W)]_{r + s - \dim(Y)}
$$
Comme $f_*[V] = [f_*\mathcal{O}_V]_r$ d'après le
lemme \ref{lemma-push-coherent}, on a
$$
[f_*V] \cdot [W] =
\sum (-1)^i
[\text{Tor}_i^{\mathcal{O}_X}(f_*\mathcal{O}_V,
\mathcal{O}_W)]_{r + s - \dim(Y)}
$$
à nouveau d'après le lemme \ref{lemma-tor-sheaf}.
En comparant la formule pour $f_*([V] \cdot f^*[W])$ à
la formule pour $f_*[V] \cdot [W]$ et en considérant le
coefficient de $Z$ par les longueurs des germes en $\xi$, on voit que
(\ref{equation-stalks}) achève la démonstration.
\end{proof}

\begin{lemma}
\label{lemma-transfer}
Soit $X \to P$ une immersion fermée de variétés non singulières.
Soit $C' \subset P \times \mathbf{P}^1$ une sous-variété fermée de dimension
$r + 1$. Supposons que
\begin{enumerate}
\item la fibre $C = C'_0$ soit de dimension $r$, c'est-à-dire que $C' \to \mathbf{P}^1$
soit dominant,
\item $C'$ coupe $X \times \mathbf{P}^1$ proprement,
\item $[C]_r$ coupe $X$ proprement.
\end{enumerate}
Alors, en posant $\alpha = [C]_r \cdot X$, considéré comme cycle sur $X$, et
$\beta = C' \cdot X \times \mathbf{P}^1$, considéré comme cycle sur
$X \times \mathbf{P}^1$, on a
$$
\alpha = \text{pr}_{X, *}(\beta \cdot X \times 0)
$$
comme cycles sur $X$, où $\text{pr}_X : X \times \mathbf{P}^1 \to X$ est la
projection.
\end{lemma}

\begin{proof}
Soit $\text{pr} : P \times \mathbf{P}^1 \to P$ la projection.
Puisque nous démontrons une égalité de cycles, il suffit de considérer
$\alpha$, resp.\ $\beta$, comme un cycle sur $P$, resp.\ $P \times \mathbf{P}^1$,
et de démontrer le résultat après image directe par $\text{pr}$.
Comme $\text{pr}^*X = X \times \mathbf{P}^1$ et que
$\text{pr}$ induit un isomorphisme de $C'_0$ sur $C$,
la formule de projection (lemme \ref{lemma-projection-formula-flat})
donne
$$
\text{pr}_*([C'_0]_r \cdot X \times \mathbf{P}^1) = [C]_r \cdot X = \alpha
$$
D'autre part, on a $[C'_0]_r = C' \cdot P \times 0$
comme cycles sur $P \times \mathbf{P}^1$,
d'après le lemme \ref{lemma-rational-equivalence-and-intersection}.
Par conséquent,
$$
[C'_0]_r \cdot X \times \mathbf{P}^1 =
(C' \cdot P \times 0) \cdot X \times \mathbf{P}^1 =
(C' \cdot X \times \mathbf{P}^1) \cdot P \times 0
$$
par associativité (lemme \ref{lemma-associative}) et commutativité du
produit d'intersection. Il reste à montrer que le produit d'intersection de
$C' \cdot X \times \mathbf{P}^1$ avec $P \times 0$ sur
$P \times \mathbf{P}^1$ est égal, en tant que cycle, au produit d'intersection de
$\beta$ avec $X \times 0$ sur $X \times \mathbf{P}^1$. Écrivons
$C' \cdot X \times \mathbf{P}^1 = \sum n_k[E_k]$, et donc
$\beta = \sum n_k[E_k]$, pour certaines sous-variétés
$E_k \subset X \times \mathbf{P}^1 \subset P \times \mathbf{P}^1$.
Les deux intersections sont alors égales à $\sum m_k[E_{k, 0}]$ d'après le
lemme \ref{lemma-rational-equivalence-and-intersection}, appliqué deux fois.
Cela achève la démonstration.
\end{proof}




\section{Projections}
\label{section-projection}

\noindent
Rappelons que nous travaillons sur un corps de base algébriquement clos fixé
$\mathbf{C}$. Si $V$ est un espace vectoriel de dimension finie sur $\mathbf{C}$,
nous posons
$$
\mathbf{P}(V) = \text{Proj}(\text{Sym}(V))
$$
où $\text{Sym}(V)$ est l'algèbre symétrique de $V$ sur $\mathbf{C}$.
Voir Constructions, exemple \ref{constructions-example-projective-space}.
La convention est choisie de telle sorte que
$V = \Gamma(\mathbf{P}(V), \mathcal{O}_{\mathbf{P}(V)}(1))$.
Bien entendu, on a $\mathbf{P}(V) \cong \mathbf{P}^n_{\mathbf{C}}$ si
$\dim(V) = n + 1$. Notons que $\mathbf{P}(V)$ est une variété projective
non singulière.

\medskip\noindent
Soit $p \in \mathbf{P}(V)$ un point fermé. Le point $p$ correspond à une
surjection $V \to L_p$ d'espaces vectoriels avec $\dim(L_p) = 1$, voir
Constructions, lemme \ref{constructions-lemma-apply}.
Notons $W_p = \Ker(V \to L_p)$.
La {\it projection de centre $p$} est le morphisme
$$
r_p : \mathbf{P}(V) \setminus \{p\} \longrightarrow \mathbf{P}(W_p)
$$
de Constructions, lemme \ref{constructions-lemma-morphism-proj}.
Voici un lemme préliminaire.

\begin{lemma}
\label{lemma-projection-generically-finite}
Soit $V$ un espace vectoriel de dimension $n + 1$.
Soit $X \subset \mathbf{P}(V)$ un sous-schéma fermé.
Si $X \not = \mathbf{P}(V)$, il existe un ouvert de Zariski non vide
$U \subset \mathbf{P}(V)$
tel que, pour tout point fermé $p \in U$, la restriction
de la projection $r_p$ définisse un morphisme fini
$r_p|_X : X \to \mathbf{P}(W_p)$.
\end{lemma}

\begin{proof}
Le lemme vaut, affirmons-nous, avec $U = \mathbf{P}(V) \setminus X$. Pour tout point
fermé $p$ de $U$, on obtient en effet un morphisme $r_p|_X : X \to \mathbf{P}(W_p)$.
Ce morphisme est propre, car $X$ est un schéma propre
(Morphismes, lemmes \ref{morphisms-lemma-locally-projective-proper} et
\ref{morphisms-lemma-image-proper-scheme-closed}). D'autre part, les
fibres de $r_p$ sont des droites affines, comme le montre un calcul direct.
Ainsi, les fibres de $r_p|X$ sont propres et affines, donc finies
(Morphismes, lemme \ref{morphisms-lemma-finite-proper}).
Enfin, un morphisme propre à fibres finies est fini
(Cohomologie des schémas, lemme \ref{coherent-lemma-characterize-finite}).
\end{proof}

\begin{lemma}
\label{lemma-projection-generically-immersion}
Soit $V$ un espace vectoriel de dimension $n + 1$.
Soit $X \subset \mathbf{P}(V)$ une sous-variété fermée.
Soit $x \in X$ un point non singulier.
\begin{enumerate}
\item Si $\dim(X) < n - 1$, il existe un ouvert de Zariski non vide
$U \subset \mathbf{P}(V)$ tel que, pour tout point fermé $p \in U$, le
morphisme $r_p|_X : X \to r_p(X)$ soit un
isomorphisme au-dessus d'un voisinage ouvert de $r_p(x)$.
\item Si $\dim(X) = n - 1$, il existe un ouvert de Zariski non vide
$U \subset \mathbf{P}(V)$ tel que, pour tout point fermé $p \in U$, le
morphisme $r_p|_X : X \to \mathbf{P}(W_p)$ soit \'etale en $x$.
\end{enumerate}
\end{lemma}

\begin{proof}
Démonstration de (1). Notons que si $x, y \in X$ ont la même image par
$r_p$, alors $p$ appartient à la droite $\overline{xy}$.
Considérons le schéma de type fini
$$
T = \{(y, p) \mid
y \in X \setminus \{x\},\ p \in \mathbf{P}(V),\ p \in \overline{xy}\}
$$
et les morphismes $T \to X$ et $T \to \mathbf{P}(V)$ donnés par
$(y, p) \mapsto y$ et $(y, p) \mapsto p$.
Puisque toute fibre de $T \to X$ est une droite, on voit que
la dimension de $T$ est $\dim(X) + 1 < \dim(\mathbf{P}(V))$.
Ainsi, $T \to \mathbf{P}(V)$ n'est pas surjectif. D'autre part,
considérons le schéma de type fini
$$
T' = \{p \mid
p \in \mathbf{P}(V) \setminus \{x\},
\ \overline{xp}\text{ tangent to }X\text{ at }x\}
$$
La dimension de $T'$ est alors $\dim(X) < \dim(\mathbf{P}(V))$.
Ainsi, le morphisme $T' \to \mathbf{P}(V)$ n'est pas non plus surjectif.
Soit $U \subset \mathbf{P}(V) \setminus X$ un ouvert non vide disjoint
de ces images ; un tel $U$ existe, car les images de $T$ et de $T'$
dans $\mathbf{P}(V)$ sont constructibles d'après
Morphismes, lemme \ref{morphisms-lemma-chevalley}.
Alors, pour $p \in U$ fermé, la projection
$r_p|_X : X \to \mathbf{P}(W_p)$ induit une application injective sur
l'espace tangent en $x$ et $r_p^{-1}(\{r_p(x)\}) = \{x\}$.
Cela signifie que $r_p$ est non ramifié en $x$
(Variétés, lemme \ref{varieties-lemma-injective-tangent-spaces-unramified}),
fini d'après le lemme \ref{lemma-projection-generically-finite},
et que $r_p^{-1}(\{r_p(x)\}) = \{x\}$ ; le lemme de Morphismes \'etales
\ref{etale-lemma-finite-unramified-one-point} s'applique donc, et
il existe un voisinage ouvert $R$ de $r_p(x)$
dans $\mathbf{P}(W_p)$ tel que $(r_p|_X)^{-1}(R) \to R$ soit une
immersion fermée, ce qui démontre (1).

\medskip\noindent
Démonstration de (2). Dans ce cas, on conclut encore que le morphisme
$T' \to \mathbf{P}(V)$ n'est pas surjectif.
En raisonnant comme ci-dessus, on conclut que, pour $U \subset \mathbf{P}(V)$
évitant $X$ et l'image de $T'$, la projection
$r_p|_X : X \to \mathbf{P}(W_p)$ est \'etale en $x$ et finie.
\end{proof}

\begin{lemma}
\label{lemma-projection-generically-immersion-refined}
Soit $V$ un espace vectoriel. Soient $Z \subset X \subset \mathbf{P}(V)$ des
sous-variétés fermées avec $Z \not = X \not = \mathbf{P}(V)$. Soit $x \in Z$ un point
non singulier sur $X$. Il existe alors un ouvert de Zariski non vide
$U \subset \mathbf{P}(V)$ tel que, pour tout point fermé $p \in U$, le
morphisme $r_p|_Z : Z \to r_p(Z)$ soit un isomorphisme au-dessus d'un voisinage ouvert
de $r_p(x)$.
\end{lemma}

\begin{proof}
En raisonnant comme dans la démonstration du lemme \ref{lemma-projection-generically-immersion},
on trouve un ouvert de Zariski non vide $U$ tel que, pour $p \in U$, l'application
$r_p|_X : X \to r_p(X)$ soit \'etale en $x$ et que
$r_p^{-1}(r_p(\{x\})) \cap Z = \{x\}$. Détails omis.
Posons $Z' = r_p(Z)$, considéré comme sous-variété fermée de $\mathbf{P}(W_p)$.
Alors le morphisme $\pi : Z \to Z'$ est fini
(lemme \ref{lemma-projection-generically-finite}), non ramifié en $x$
(petit détail omis) et $\pi^{-1}(\pi(\{x\})) = \{x\}$. Le résultat découle
du lemme de Morphismes \'etales \ref{etale-lemma-finite-unramified-one-point}.
\end{proof}

\begin{lemma}
\label{lemma-projection-injective}
Soit $V$ un espace vectoriel de dimension $n + 1$.
Soient $Y, Z \subset \mathbf{P}(V)$ des sous-variétés fermées.
Il existe un ouvert de Zariski non vide $U \subset \mathbf{P}(V)$
tel que, pour tout point fermé $p \in U$, on ait
$$
Y \cap r_p^{-1}(r_p(Z)) = (Y \cap Z) \cup E
$$
où $E \subset Y$ est fermé et
$\dim(E) \leq \dim(Y) + \dim(Z) + 1 - n$.
\end{lemma}

\begin{proof}
Posons $Y' = Y \setminus Y \cap Z$.
Soient $y \in Y'$, $z \in Z$ des points fermés tels que $r_p(y) = r_p(z)$.
Alors $p$ appartient à la droite $\overline{yz}$ passant par $y$ et $z$.
Considérons le schéma de type fini
$$
T = \{(y, z, p) \mid y \in Y', z \in Z, p \in \overline{yz}\}
$$
et le morphisme $T \to \mathbf{P}(V)$ donné par $(y, z, p) \mapsto p$.
Observons que $T$ est irréductible et que $\dim(T) = \dim(Y) + \dim(Z) + 1$.
La fibre générale de $T \to \mathbf{P}(V)$ est donc de dimension au plus
$\dim(Y) + \dim(Z) + 1 - n$ ; plus précisément, il existe un ouvert non vide
$U \subset \mathbf{P}(V) \setminus (Y \cup Z)$ au-dessus
duquel la fibre est de dimension au plus $\dim(Y) + \dim(Z) + 1 - n$
(Variétés, lemme \ref{varieties-lemma-dimension-fibres-locally-algebraic}).
Soit $p \in U$ un point fermé et soit $F \subset T$ la fibre
de $T \to \mathbf{P}(V)$ au-dessus de $p$. Alors
$$
(Y \cap r_p^{-1}(r_p(Z))) \setminus (Y \cap Z)
$$
est l'image de $F \to Y$, $(y, z, p) \mapsto y$. À nouveau d'après
Variétés, lemme \ref{varieties-lemma-dimension-fibres-locally-algebraic},
l'adhérence de l'image de $F \to Y$ est de dimension au plus
$\dim(Y) + \dim(Z) + 1 - n$.
\end{proof}

\begin{lemma}
\label{lemma-find-lines}
Soit $V$ un espace vectoriel. Soit $B \subset \mathbf{P}(V)$
une sous-variété fermée de codimension $\geq 2$.
Soit $p \in \mathbf{P}(V)$ un point fermé, $p \not \in B$.
Il existe alors une droite $\ell \subset \mathbf{P}(V)$
avec $p \in \ell$ et $\ell \cap B = \emptyset$. De plus, ces droites
recouvrent un ouvert de $\mathbf{P}(V)$.
\end{lemma}

\begin{proof}
Considérons l'image de $B$ par la projection
$r_p : \mathbf{P}(V) \to \mathbf{P}(W_p)$.
Puisque $\dim(W_p) = \dim(V) - 1$, on voit que $r_p(B)$
est de codimension $\geq 1$ dans $\mathbf{P}(W_p)$.
Pour tout $q \in \mathbf{P}(V)$ tel que $r_p(q) \not \in r_p(B)$,
on voit que la droite $\ell = \overline{pq}$ joignant $p$ et $q$ convient.
\end{proof}

\begin{lemma}
\label{lemma-doubly-transitive}
Soit $V$ un espace vectoriel. Soit $G = \text{PGL}(V)$.
Alors $G \times \mathbf{P}(V) \to \mathbf{P}(V)$ est
doublement transitif.
\end{lemma}

\begin{proof}
Omis. Indication : cela résulte du fait que $\text{GL}(V)$ agit de façon
doublement transitive sur les couples de vecteurs linéairement indépendants.
\end{proof}

\begin{lemma}
\label{lemma-determinant}
Soit $k$ un corps. Soit $n \geq 1$ un entier et soient
$x_{ij}, 1 \leq i, j \leq n$ des indéterminées. Alors
$$
\det
\left(
\begin{matrix}
x_{11} & x_{12} & \ldots & x_{1n} \\
x_{21} & \ldots & \ldots & \ldots \\
\ldots & \ldots & \ldots & \ldots \\
x_{n1} & \ldots & \ldots & x_{nn}
\end{matrix}
\right)
$$
est un élément irréductible de l'anneau de polynômes $k[x_{ij}]$.
\end{lemma}

\begin{proof}
Soit $V$ un espace vectoriel de dimension $n$. Géométriquement,
le lemme signifie que le lieu $C$ des applications linéaires non inversibles
$V \to V$ est irréductible (notons que, comme le déterminant est de degré
$1$ en chaque indéterminée, il est sans facteur carré). Soit $W$ un espace vectoriel de
dimension $n - 1$. Par algèbre linéaire élémentaire, le morphisme
$$
\Hom(W, V) \times \Hom(V, W) \longrightarrow \Hom(V, V),\quad
(\psi, \varphi) \longmapsto \psi \circ \varphi
$$
a pour image $C$. Puisque la source est irréductible, l'image l'est aussi.
Détails omis.
\end{proof}

\noindent
Soit $V$ un espace vectoriel de dimension $n + 1$. Posons $E = \text{End}(V)$.
Soit $E^\vee = \Hom(E, \mathbf{C})$ l'espace vectoriel dual.
Écrivons $\mathbf{P} = \mathbf{P}(E^\vee)$.
Il existe une application linéaire canonique
$$
V \longrightarrow V \otimes_\mathbf{C} E^\vee = \Hom(E, V)
$$
qui envoie $v \in V$ sur l'application $g \mapsto g(v)$ de $\Hom(E, V)$.
Rappelons que nous avons une application canonique
$E^\vee \to \Gamma(\mathbf{P}, \mathcal{O}_\mathbf{P}(1))$
qui est un isomorphisme. On obtient donc une application canonique
$$
\psi : V \otimes \mathcal{O}_\mathbf{P} \to V \otimes \mathcal{O}_\mathbf{P}(1)
$$
de faisceaux de modules sur $\mathbf{P}$ qui, sur les sections globales, redonne
l'application donnée. Rappelons qu'un fibré projectif $\mathbf{P}(\mathcal{E})$
est défini comme le Proj relatif de l'algèbre symétrique de $\mathcal{E}$, voir
Constructions, définition \ref{constructions-definition-projective-bundle}.
Nous allons étudier l'application rationnelle
entre $\mathbf{P}(V \otimes \mathcal{O}_\mathbf{P}(1))$ et
$\mathbf{P}(V \otimes \mathcal{O}_\mathbf{P})$ associée à $\psi$. D'après
Constructions, lemme \ref{constructions-lemma-relative-proj-base-change},
nous avons un isomorphisme canonique
$$
\mathbf{P}(V \otimes \mathcal{O}_\mathbf{P}) = \mathbf{P} \times \mathbf{P}(V)
$$
D'après Constructions, lemme \ref{constructions-lemma-twisting-and-proj},
on voit que
$$
\mathbf{P}(V \otimes \mathcal{O}_\mathbf{P}(1)) =
\mathbf{P}(V \otimes \mathcal{O}_\mathbf{P}) = \mathbf{P} \times \mathbf{P}(V)
$$
En combinant ceci avec
Constructions, lemme \ref{constructions-lemma-morphism-relative-proj},
on obtient
\begin{equation}
\label{equation-r-psi}
\mathbf{P} \times \mathbf{P}(V) \supset
U(\psi) \xrightarrow{r_\psi} \mathbf{P} \times \mathbf{P}(V)
\end{equation}
Pour mieux comprendre ceci, étudions ce qui se passe sur les fibres au-dessus de
$\mathbf{P}$. Soit $g \in E$ non nul. Il définit une application non nulle
$E^\vee \to \mathbf{C}$, donc un point $[g] \in \mathbf{P}$.
D'autre part, $g$ définit une application $\mathbf{C}$-linéaire $g : V \to V$.
Ainsi, d'après
Constructions, lemme \ref{constructions-lemma-morphism-proj},
on obtient une application
$$
\mathbf{P}(V) \supset U(g) \xrightarrow{r_g} \mathbf{P}(V)
$$
Nous utiliserons ci-dessous le fait que $U(g)$ est la fibre $U(\psi)_{[g]}$ et
que $r_g$ est la fibre de $r_\psi$ au-dessus du point $[g]$. Une autre observation
que nous utiliserons est que le complémentaire de $U(g)$ dans $\mathbf{P}(V)$ est
l'image de l'immersion fermée
$$
\mathbf{P}(\Coker(g)) \longrightarrow \mathbf{P}(V)
$$
et que l'image de $r_g$ est l'image de l'immersion fermée
$$
\mathbf{P}(\Im(g)) \longrightarrow \mathbf{P}(V)
$$

\begin{lemma}
\label{lemma-make-family}
Avec les notations ci-dessus. Soient $X, Y$ des sous-variétés fermées de $\mathbf{P}(V)$
qui se coupent proprement, telles que $X \not = \mathbf{P}(V)$ et
$X \cap Y \not = \emptyset$. Pour une droite générale $\ell \subset \mathbf{P}$
avec $[\text{id}_V] \in \ell$, on a
\begin{enumerate}
\item $X \subset U_g$ pour tout $[g] \in \ell$,
\item $g(X)$ coupe $Y$ proprement pour tout $[g] \in \ell$.
\end{enumerate}
\end{lemma}

\begin{proof}
Soit $B \subset \mathbf{P}$ l'ensemble des points «~mauvais~», c'est-à-dire des
points $[g]$ qui ne vérifient pas (1) ou (2). Notons que
$[\text{id}_V] \not \in B$ par hypothèse. De plus, $B$ est fermé.
Il suffit donc de montrer que $\dim(B) \leq \dim(\mathbf{P}) - 2$
(lemme \ref{lemma-find-lines}).

\medskip\noindent
Considérons d'abord l'ouvert $G = \text{PGL}(V) \subset \mathbf{P}$
formé des points $[g]$ tels que $g : V \to V$ soit inversible.
Puisque $G$ agit doublement transitivement sur $\mathbf{P}(V)$
(lemme \ref{lemma-doubly-transitive}),
on voit que
$$
T = \{(x, y, [g]) \mid x \in X, y \in Y, [g] \in G, r_g(x) = y\}
$$
est une fibration localement triviale sur $X \times Y$ dont la fibre est égale
au stabilisateur d'un point dans $G$. Ainsi, $T$ est une variété.
Observons que la fibre de $T \to G$ au-dessus de $[g]$ est $r_g(X) \cap Y$.
Le morphisme $T \to G$ est surjectif, car toute transformée de $X$
rencontre $Y$ (notons que l'hypothèse selon laquelle $X$ et $Y$ se coupent
proprement et que $X \cap Y \not = \emptyset$ entraîne
$\dim(X) + \dim(Y) \geq \dim(\mathbf{P}(V))$, puis le
lemme de Variétés \ref{varieties-lemma-intersection-in-projective-space}
implique que toute transformée de $X$ rencontre $Y$).
Puisque les dimensions des fibres d'un morphisme dominant ne
sautent pas en codimension $1$
(Variétés, lemme \ref{varieties-lemma-dimension-fibres-locally-algebraic}),
on conclut que $B \cap G$ est de codimension $\geq 2$.

\medskip\noindent
Considérons ensuite le complémentaire $Z = \mathbf{P} \setminus G$.
C'est une variété irréductible, car le déterminant est un
polynôme irréductible (lemme \ref{lemma-determinant}).
Il suffit donc de montrer que $B$ ne contient pas le
point générique de $Z$. Pour un point général $[g] \in Z$,
le conoyau $V \to \Coker(g)$ est de dimension $1$, donc
$U(g)$ est le complémentaire d'un point. Comme $X \not = \mathbf{P}(V)$,
on voit que, pour $[g] \in Z$ général, on a $X \subset U(g)$.
De plus, le morphisme $r_g|_X : X \to r_g(X)$ est fini, donc
$\dim(r_g(X)) = \dim(X)$.
D'autre part, pour un tel $g$, l'image de $r_g$ est le
sous-espace fermé $H = \mathbf{P}(\Im(g)) \subset \mathbf{P}(V)$
qui est de codimension $1$.
Pour un point général de $Z$, on voit que $H \cap Y$ est de dimension $1$
de moins que $Y$ (comparer à
Variétés, lemme \ref{varieties-lemma-exact-sequence-induction}).
On voit ainsi qu'il faut montrer que $r_g(X)$ et $H \cap Y$
se coupent proprement dans $H$. Pour un choix fixé de $H$, on peut,
en composant $g$ à gauche avec un automorphisme, déplacer $r_g(X)$ par
un automorphisme arbitraire de $H = \mathbf{P}(\Im(g))$.
On peut donc raisonner comme ci-dessus pour conclure que l'intersection
de $H \cap Y$ avec $r_g(X)$ est propre pour $g$ général ayant
$H = \mathbf{P}(\Im(g))$. Certains détails sont omis.
\end{proof}







\section{Lemme de déplacement}
\label{section-moving-lemma}

\noindent
Le lemme de déplacement affirme qu'étant donnés un $r$-cycle $\alpha$ et un $s$-cycle
$\beta$, il existe $\alpha'$, $\alpha' \sim_{rat} \alpha$, tel que
$\alpha'$ et $\beta$ se coupent proprement (lemme \ref{lemma-moving-move}).
Voir \cite{Samuel}, \cite{ChevalleyI}, \cite{ChevalleyII}.
Le point clé est le lemme \ref{lemma-moving} ; le lecteur pourra trouver
ce lemme sous la forme énoncée dans
\cite[Example 11.4.1]{F} et une démonstration dans \cite{Roberts}.

\begin{lemma}
\label{lemma-moving}
\begin{reference}
Voir \cite{Roberts}.
\end{reference}
Soit $X \subset \mathbf{P}^N$ une sous-variété fermée non singulière.
Posons $n = \dim(X)$ et $0 \leq d, d' < n$. Soit $Z \subset X$ une sous-variété
fermée de dimension $d$ et soient $T_i \subset X$, $i \in I$, une famille finie
de sous-variétés fermées de dimension $d'$. Il existe alors
une sous-variété $C \subset \mathbf{P}^N$ telle que $C$ coupe $X$
proprement et tel que
$$
C \cdot X = Z + \sum\nolimits_{j \in J} m_j Z_j
$$
où les $Z_j \subset X$ sont irréductibles de dimension $d$, distinctes de $Z$, et
$$
\dim(Z_j \cap T_i) \leq \dim(Z \cap T_i)
$$
avec inégalité stricte si $Z$ ne rencontre pas $T_i$ proprement dans $X$.
\end{lemma}

\begin{proof}
Écrivons $\mathbf{P}^N = \mathbf{P}(V_N)$, de sorte que $\dim(V_N) = N + 1$.
Nous allons choisir une suite de projections depuis des points
\begin{align*}
& r_N : 
\mathbf{P}(V_N) \setminus \{p_N\} \to \mathbf{P}(V_{N - 1}), \\
& r_{N - 1} :
\mathbf{P}(V_{N - 1}) \setminus \{p_{N - 1}\} \to \mathbf{P}(V_{N - 2}), \\
& \ldots,\\
& r_{n + 1} :
\mathbf{P}(V_{n + 1}) \setminus \{p_{n + 1}\} \to \mathbf{P}(V_n)
\end{align*}
comme à la Section \ref{section-projection}. À chaque étape, nous choisirons
$p_N, p_{N - 1}, \ldots, p_{n + 1}$ dans un ouvert de Zariski convenable.
Choisissons un point fermé $x \in Z \subset X$. Pour tout $i$, choisissons des
points fermés $x_{it} \in T_i \cap Z$, au moins un dans chaque composante irréductible
de $T_i \cap Z$. En prenant la composée, nous obtenons
un morphisme
$$
\pi = (r_{n + 1} \circ \ldots \circ r_N)|_X :
X \longrightarrow \mathbf{P}(V_n)
$$
qui possède les propriétés suivantes :
\begin{enumerate}
\item $\pi$ est fini,
\item $\pi$ est \'etale en $x$ et en tous les $x_{it}$,
\item $\pi|_Z : Z \to \pi(Z)$ est un isomorphisme
au-dessus d'un voisinage ouvert de $\pi(x_{it})$,
\item $T_i \cap \pi^{-1}(\pi(Z)) = (T_i \cap Z) \cup E_i$, où
$E_i \subset T_i$ est fermé et
$\dim(E_i) \leq d + d' + 1 - (n + 1) = d + d' - n$.
\end{enumerate}
Il résulte immédiatement des
Lemmes \ref{lemma-projection-generically-finite},
\ref{lemma-projection-generically-immersion},
\ref{lemma-projection-generically-immersion-refined} et
\ref{lemma-projection-injective}, ainsi que d'une récurrence, que ce choix est possible ;
observons que la dernière projection part de $\mathbf{P}(V_{n + 1})$ et que
$\dim(V_{n + 1}) = n + 2$, ce qui explique l'inégalité de (4). Plus
précisément, soient $X_j, Z_j, T_{i, j} \subset \mathbf{P}(V_j)$ les images
de $X$, $Z$, $T_i$ pour $j = N, \ldots, n$. Alors
$T_{i, j + 1} \cap r_{j + 1}^{-1}(Z_j) = (T_{i, j + 1} \cap Z_{j + 1})
\cup E_{i, j + 1}$, où nous disposons d'une borne sur la dimension de $E_{i, j + 1}$
d'après le Lemme \ref{lemma-projection-injective}. Notons alors que
$T_i \cap \pi^{-1}(\pi(Z))$ est la réunion de $T_i \cap Z$
et des images réciproques des $E_{i, j}$. Comme les morphismes de transition
entre les $X_j$ sont finis, on obtient la borne de dimension voulue.

\medskip\noindent
Soit $C \subset \mathbf{P}(V_N)$ l'adhérence schématique de
$(r_{n + 1} \circ \ldots \circ r_N)^{-1}(\pi(Z))$. Puisque $\pi$
est \'etale au point $x$ de $Z$, le sous-schéma fermé
$C \cap X = \pi^{-1}(\pi(Z))$
contient $Z$ avec multiplicité $1$ (le calcul local est omis).
Le Lemme \ref{lemma-transversal-subschemes} donne donc
$$
C \cdot X = [Z] + \sum m_j[Z_j]
$$
pour certaines sous-variétés $Z_j \subset X$ de dimension $d$. Notons que
$$
C \cap X = \pi^{-1}(\pi(Z))
$$
ensemblistement. Par conséquent,
$T_i \cap Z_j \subset T_i \cap \pi^{-1}(\pi(Z)) \subset (T_i \cap Z) \cup E_i$.
Pour toute composante irréductible de $T_i \cap Z_j$ contenue dans $E_i$, nous
avons la borne de dimension voulue. Enfin, soit $V$ une composante irréductible
de $T_i \cap Z_j$ contenue dans $T_i \cap Z$. Pour achever
la démonstration, il suffit de montrer que $V$ ne contient aucun des
points $x_{it}$, car alors $\dim(V) < \dim(Z \cap T_i)$.
Pour cela, il suffit de montrer que $x_{it} \not \in Z_j$
pour tous $i, t, j$.

\medskip\noindent
Posons $Z' = \pi(Z)$ et $Z'' = \pi^{-1}(Z')$, schématiquement. D'après
la condition (3), il existe un ouvert $U \subset \mathbf{P}(V_n)$ contenant
$\pi(x_{it})$ tel que $\pi^{-1}(U) \cap Z \to U \cap Z'$ soit un isomorphisme.
En particulier, $Z \to Z'$ est un isomorphisme local en $x_{it}$.
D'autre part, $Z'' \to Z'$ est \'etale en $x_{it}$ d'après la condition (2).
L'immersion fermée $Z \to Z''$ est donc \'etale en $x_{it}$
(Morphismes, Lemme \ref{morphisms-lemma-etale-permanence}).
Ainsi $Z = Z''$ dans un voisinage de Zariski de $x_{it}$, ce qui démontre
l'assertion.
\end{proof}

\noindent
Le déplacement proprement dit s'effectue à l'aide du lemme suivant.

\begin{lemma}
\label{lemma-move}
Soit $C \subset \mathbf{P}^N$ une sous-variété fermée.
Soit $X \subset \mathbf{P}^N$ une sous-variété et soient $T_i \subset X$
une collection finie de sous-variétés fermées.
Supposons que $C$ et $X$ se rencontrent proprement.
Alors il existe une sous-variété fermée
$C' \subset \mathbf{P}^N \times \mathbf{P}^1$ telle que
\begin{enumerate}
\item $C' \to \mathbf{P}^1$ soit dominant,
\item $C'_0 = C$ schématiquement,
\item $C'$ et $X \times \mathbf{P}^1$ se rencontrent proprement,
\item $C'_\infty$ rencontre proprement chacun des $T_i$ donnés.
\end{enumerate}
\end{lemma}

\begin{proof}
Si $C \cap X = \emptyset$, prenons la famille constante
$C' = C \times \mathbf{P}^1$. Nous pouvons donc supposer
$C \cap X \not = \emptyset$.

\medskip\noindent
Écrivons $\mathbf{P}^N = \mathbf{P}(V)$, de sorte que $\dim(V) = N + 1$. Soit
$E = \text{End}(V)$. Soit $E^\vee = \Hom(E, \mathbf{C})$. Posons
$\mathbf{P} = \mathbf{P}(E^\vee)$ comme dans le Lemme \ref{lemma-make-family}.
Choisissons une droite générale $\ell \subset \mathbf{P}$ passant par $\text{id}_V$.
Définissons $C' \subset \ell \times \mathbf{P}(V)$ comme le
sous-schéma fermé de fibre $r_g(C)$ au-dessus de $[g] \in \ell$.
Plus précisément, $C'$ est l'image de
$$
\ell \times C \subset \mathbf{P} \times \mathbf{P}(V)
$$
par le morphisme (\ref{equation-r-psi}). D'après le Lemme \ref{lemma-make-family},
cela a un sens, c'est-à-dire que $\ell \times C \subset U(\psi)$. Le morphisme
$\ell \times C \to C'$ est fini et $C'_{[g]} = r_g(C)$ ensemblistement
pour tout $[g] \in \ell$. Les parties (1) et (2) sont claires avec
$0 = [\text{id}_V] \in \ell$. La partie (3) résulte du fait
que $r_g(C)$ et $X$ se rencontrent proprement pour tout $[g] \in \ell$.
La partie (4) résulte du fait qu'un point général $\infty = [g] \in \ell$
est un point général de $\mathbf{P}$ et que, pour un tel point,
$r_g(C) \cap T$ est propre pour toute sous-variété fermée $T$ de $\mathbf{P}(V)$.
Les détails sont omis.
\end{proof}

\begin{lemma}
\label{lemma-moving-move}
Soit $X$ une variété projective non singulière. Soient $\alpha$ un
$r$-cycle et $\beta$ un $s$-cycle sur $X$. Alors il existe
un $r$-cycle $\alpha'$ tel que $\alpha' \sim_{rat} \alpha$ et
tel que $\alpha'$ et $\beta$ se rencontrent proprement.
\end{lemma}

\begin{proof}
Écrivons $\beta = \sum n_i[T_i]$ pour certaines sous-variétés $T_i \subset X$
de dimension $s$. Par linéarité, nous pouvons supposer que $\alpha = [Z]$ pour
une sous-variété fermée irréductible $Z \subset X$ de dimension $r$.
Nous démontrerons le lemme par récurrence sur le maximum $e$ des entiers
$$
\dim(Z \cap T_i)
$$
Le cas initial est $e = r + s - \dim(X)$. Dans ce cas, $Z$ rencontre
$\beta$ proprement et le lemme est trivial.

\medskip\noindent
Étape de récurrence. Supposons $e > r + s - \dim(X)$. Choisissons un plongement
$X \subset \mathbf{P}^N$ et appliquons le Lemme \ref{lemma-moving} afin de trouver une
sous-variété fermée $C \subset \mathbf{P}^N$ telle que
$C \cdot X = [Z] + \sum m_j[Z_j]$ et telle que l'hypothèse de
récurrence s'applique à chaque $Z_j$. Appliquons ensuite le Lemme \ref{lemma-move}
à $C$, $X$, $T_i$ pour trouver $C' \subset \mathbf{P}^N \times \mathbf{P}^1$.
Soit $\gamma = C' \cdot X \times \mathbf{P}^1$, considéré comme un cycle
sur $X \times \mathbf{P}^1$. D'après le Lemme \ref{lemma-transfer}, nous avons
$$
[Z] + \sum m_j[Z_j] = \text{pr}_{X, *}(\gamma \cdot X \times 0)
$$
D'autre part, le cycle
$\gamma_\infty = \text{pr}_{X, *}(\gamma \cdot X \times \infty)$
est supporté sur $C'_\infty \cap X$ et rencontre donc $\beta$ proprement.
Ainsi $[Z] \sim_{rat} - \sum m_j[Z_j] + \gamma_\infty$
d'après le Lemme \ref{lemma-rational-equivalence-and-intersection}. Comme, par
récurrence, chaque $[Z_j]$ est rationnellement équivalent à un cycle qui rencontre
$\beta$ proprement, cela achève la démonstration.
\end{proof}



\section{Produits d'intersection et équivalence rationnelle}
\label{section-intersections-and-rational-equivalence}

\noindent
Avec les définitions précédentes, montrons que le produit d'intersection
est bien défini modulo équivalence rationnelle. Traitons d'abord un
cas particulier.

\begin{lemma}
\label{lemma-well-defined-special-case}
Soit $X$ une variété non singulière. Soit
$W \subset X \times \mathbf{P}^1$ une sous-variété de dimension $(s + 1)$
dominant $\mathbf{P}^1$. Soient $W_a$, resp.\ $W_b$, les fibres de
$W \to \mathbf{P}^1$ au-dessus de $a$, resp.\ de $b$. Soit $V$ une sous-variété de dimension $r$
de $X$ telle que $V$ rencontre proprement $W_a$ et
$W_b$. Alors $[V] \cdot [W_a]_r \sim_{rat} [V] \cdot [W_b]_r$.
\end{lemma}

\begin{proof}
Nous avons $[W_a]_r = \text{pr}_{X,*}(W \cdot X \times a)$ et de même pour
$[W_b]_r$, voir le Lemme \ref{lemma-rational-equivalence-and-intersection}.
Nous nous ramenons donc à démontrer
$$
V \cdot \text{pr}_{X,*}( W \cdot X \times a) \sim_{rat} V \cdot
\text{pr}_{X,*}( W \cdot X\times b).
$$
En appliquant la formule de projection du
Lemme \ref{lemma-projection-formula-flat}, on obtient
$$
V \cdot \text{pr}_{X,*}( W \cdot X\times a) =
\text{pr}_{X,*}(V \times \mathbf{P}^1 \cdot (W \cdot X\times a))
$$
et de même pour $b$. Nous nous ramenons donc à démontrer
$$
\text{pr}_{X,*}(V \times \mathbf{P}^1 \cdot (W \cdot X\times a))
\sim_{rat}
\text{pr}_{X,*}(V \times \mathbf{P}^1 \cdot (W \cdot X\times b))
$$
Si $V \times \mathbf{P}^1$ et $W$ se rencontrent proprement, alors
l'associativité des multiplicités d'intersection
(Lemme \ref{lemma-associative})
donne $V \times \mathbf{P}^1 \cdot (W \cdot X\times a) =
(V \times \mathbf{P}^1 \cdot W) \cdot X \times a$,
et de même pour $b$. Nous nous ramenons donc à démontrer
$$
\text{pr}_{X,*}((V \times \mathbf{P}^1 \cdot W) \cdot X\times a)
\sim_{rat}
\text{pr}_{X,*}((V \times \mathbf{P}^1 \cdot W) \cdot X\times b)
$$
ce qui est vrai d'après le Lemme \ref{lemma-rational-equivalence-and-intersection}.

\medskip\noindent
Le raisonnement précédent ne fonctionne pas tout à fait. L'obstacle est que
nous ne savons pas si $V \times \mathbf{P}^1$ et $W$ se rencontrent proprement.
Nous savons seulement que $V$ et $W_a$, ainsi que $V$ et $W_b$, se rencontrent proprement.
Soient $Z_i$, $i \in I$, les composantes irréductibles de
$V \times \mathbf{P}^1 \cap W$. Nous savons alors que
$\dim(Z_i) \geq r + 1 + s + 1 - n - 1 = r + s + 1 - n$, où $n = \dim(X)$, voir
le Lemme \ref{lemma-intersect-in-smooth}. Comme nous avons supposé
que $V$ et $W_a$ se rencontrent proprement, on voit que
$\dim(Z_{i, a}) = r + s - n$ ou $Z_{i, a} = \emptyset$.
D'autre part, si $Z_{i, a} \not = \emptyset$, alors
$\dim(Z_{i, a}) \geq \dim(Z_i) - 1 = r + s - n$.
Il s'ensuit que $\dim(Z_i) = r + s + 1 - n$ si $Z_i$ rencontre $X \times a$
et, dans ce cas, $Z_i \to \mathbf{P}^1$ est surjectif.
Nous pouvons donc écrire $I = I' \amalg I''$, où $I'$ est l'ensemble des $i \in I$
tels que $Z_i \to \mathbf{P}^1$ soit surjectif, et $I''$ l'ensemble des
$i \in I$ tels que $Z_i$ soit au-dessus d'un point fermé $t_i \in \mathbf{P}^1$
avec $t_i \not = a$ et $t_i \not = b$. Considérons le cycle
$$
\gamma = \sum\nolimits_{i \in I'} e_i [Z_i]
$$
où nous posons
$$
e_i = \sum\nolimits_p (-1)^p
\text{long}_{\mathcal{O}_{X \times \mathbf{P}^1, Z_i}}
\text{Tor}_p^{\mathcal{O}_{X \times \mathbf{P}^1, Z_i}}(
\mathcal{O}_{V \times \mathbf{P}^1, Z_i}, \mathcal{O}_{W, Z_i})
$$
Nous montrerons que $\gamma$ peut remplacer
le produit d'intersection de $V \times \mathbf{P}^1$ et $W$.

\medskip\noindent
Nous le démontrerons à l'aide de l'associativité des produits d'intersection, exactement
comme ci-dessus. Soit $U = \mathbf{P}^1 \setminus \{t_i, i \in I''\}$.
Notons que $X \times a$ et $X \times b$ sont contenus dans $X \times U$.
Les sous-variétés
$$
V \times U,\quad W_U,\quad X \times a\quad\text{de}\quad X \times U
$$
se rencontrent deux à deux transversalement par notre choix de $U$ et, de plus,
$\dim(V \times U \cap W_U \cap X \times a) = \dim(V \cap W_a)$ a
la dimension attendue. On voit donc que
$$
V \times U \cdot (W_U \cdot X \times a) =
(V \times U \cdot W_U) \cdot X \times a
$$
comme cycles sur $X \times U$, d'après le Lemme \ref{lemma-associative}.
Par construction, $\gamma$ se restreint au cycle $V \times U \cdot W_U$
sur $X \times U$. Trivialement,
$V \times \mathbf{P}^1 \cdot (W \times X \times a)$ se restreint
à $V \times U \cdot (W_U \cdot X \times a)$ sur $X \times U$.
Par conséquent,
$$
V \times \mathbf{P}^1 \cdot (W \cdot X \times a) =
\gamma \cdot X \times a
$$
comme cycles sur $X \times \mathbf{P}^1$ (car les deux membres
sont contenus dans $X \times U$ et sont égaux après restriction
à $X \times U$ d'après ce qui précède). Comme on a le même résultat pour $b$,
on conclut que
\begin{align*}
V \cdot [W_a]
& =
\text{pr}_{X,*}(V \times \mathbf{P}^1 \cdot (W \cdot X\times a)) \\
& =
\text{pr}_{X, *}(\gamma \cdot X \times a) \\
& \sim_{rat} 
\text{pr}_{X, *}(\gamma \cdot X \times b) \\
& =
\text{pr}_{X,*}(V \times \mathbf{P}^1 \cdot (W \cdot X\times b)) \\
& =
V \cdot [W_b]
\end{align*}
La première et la dernière égalité résultent du premier paragraphe de la démonstration,
la deuxième et l'avant-dernière ont été démontrées dans le présent paragraphe, et
l'équivalence centrale est donnée par le
Lemme \ref{lemma-rational-equivalence-and-intersection}.
\end{proof}

\begin{theorem}
\label{theorem-well-defined}
Soit $X$ une variété projective non singulière. Soient $\alpha$, resp.\ $\beta$,
un $r$-cycle, resp.\ un $s$-cycle, sur $X$. Supposons que $\alpha$ et $\beta$
se rencontrent proprement, de sorte que $\alpha \cdot \beta$ soit défini. Enfin,
supposons que $\alpha \sim_{rat} 0$. Alors $\alpha \cdot \beta \sim_{rat} 0$.
\end{theorem}

\begin{proof}
Choisissons une immersion fermée $X \subset \mathbf{P}^N$.
Par linéarité, il suffit de démontrer le résultat lorsque $\beta = [Z]$ pour une
sous-variété fermée de dimension $s$, notée $Z \subset X$, qui rencontre $\alpha$
proprement. La condition $\alpha \sim_{rat} 0$ signifie qu'il
existe un nombre fini de sous-variétés fermées de dimension $(r + 1)$,
$W_i \subset X \times \mathbf{P}^1$, telles que
$$
\alpha = \sum [W_{i, a_i}]_r - [W_{i, b_i}]_r
$$
pour certains couples de points $a_i, b_i$ de $\mathbf{P}^1$.
Soient $W_{i, a_i}^t$ et $W_{i, b_i}^t$ les composantes irréductibles
de $W_{i, a_i}$ et $W_{i, b_i}$.
Nous raisonnons par récurrence sur le maximum $d$ des entiers
$$
\dim(Z \cap W_{i, a_i}^t),\quad \dim(Z \cap W_{i, b_i}^t)
$$
La principale difficulté dans la suite de la démonstration est que, bien que $Z$
rencontre $\alpha$ proprement, il se peut que $Z$ ne rencontre pas les
variétés « intermédiaires » $W_{i, a_i}^t$
et $W_{i, b_i}^t$ proprement, c'est-à-dire que l'on peut avoir $d >  r + s - \dim(X)$.

\medskip\noindent
Cas initial : $d = r + s - \dim(X)$. Dans ce cas, toutes les intersections de
$Z$ avec les $W_{i, a_i}^t$ et les $W_{i, b_i}^t$ sont propres et le
résultat voulu découle du Lemme \ref{lemma-well-defined-special-case},
car celui-ci montre que
$[Z] \cdot [W_{i, a_i}]_r \sim_{rat} [Z] \cdot [W_{i, b_i}]_r$ pour
tout $i$.

\medskip\noindent
Étape de récurrence : $d >  r + s - \dim(X)$. Appliquons le
Lemme \ref{lemma-moving} à $Z \subset X$ et à
la famille de sous-variétés $\{W_{i, a_i}^t, W_{i, b_i}^t\}$. On obtient alors une
sous-variété fermée $C \subset \mathbf{P}^N$ qui rencontre $X$
proprement et telle que
$$
C \cdot X = [Z] + \sum m_j [Z_j]
$$
et telle que
$$
\dim(Z_j \cap W_{i, a_i}^t) \leq \dim(Z \cap W_{i, a_i}^t),\quad
\dim(Z_j \cap W_{i, b_i}^t) \leq \dim(Z \cap W_{i, b_i}^t)
$$
avec inégalité stricte si le membre de droite est $> r + s - \dim(X)$.
Cela implique deux choses : (a) l'hypothèse de récurrence s'applique à
chaque $Z_j$, et (b) $C \cdot X$ et $\alpha$ se rencontrent proprement (car
$\alpha$ est une combinaison linéaire des $[W_{i, a_i}^t]$ et
$[W_{i, a_i}^t]$ qui rencontrent $Z$ proprement).
Choisissons ensuite $C' \subset \mathbf{P}^N \times \mathbf{P}^1$
comme dans le Lemme \ref{lemma-move}, relativement à $C$, $X$ et
$W_{i, a_i}^t$, $W_{i, b_i}^t$.
Écrivons $C' \cdot X \times \mathbf{P}^1 = \sum n_k [E_k]$ pour
certaines sous-variétés $E_k \subset X \times \mathbf{P}^1$ de
dimension $s + 1$. Notons que $n_k > 0$ pour tout $k$, d'après la
Proposition \ref{proposition-positivity}.
D'après le Lemme \ref{lemma-transfer}, nous avons
$$
[Z] + \sum m_j [Z_j] = \sum n_k[E_{k, 0}]_s
$$
Comme $E_{k, 0} \subset C \cap X$, on voit que $[E_{k, 0}]_s$ et $\alpha$
se rencontrent proprement. D'autre part, le cycle
$$
\gamma = \sum n_k[E_{k, \infty}]_s
$$
est supporté sur $C'_\infty \cap X$ et
rencontre donc proprement chaque $W_{i, a_i}^t$, $W_{i, b_i}^t$.
Par le cas initial et la linéarité, on voit donc que
$$
\gamma \cdot \alpha \sim_{rat} 0
$$
Comme nous avons vu que $E_{k, 0}$ et $E_{k, \infty}$
rencontrent $\alpha$ proprement, le Lemme \ref{lemma-well-defined-special-case},
appliqué à $E_k \subset X \times \mathbf{P}^1$ et $\alpha$, donne
$$
[E_{k, 0}] \cdot \alpha \sim_{rat} [E_{k, \infty}] \cdot \alpha
$$
En regroupant le tout, nous avons
\begin{align*}
[Z] \cdot \alpha
& =
(\sum n_k[E_{k, 0}]_r - \sum m_j[Z_j]) \cdot \alpha \\
& \sim_{rat}
\sum n_k [E_{k, 0}] \cdot \alpha \quad (\text{par l'hypothèse de récurrence})\\
& \sim_{rat}
\sum n_k [E_{k, \infty}] \cdot \alpha \quad (\text{par le lemme})\\
& =
\gamma \cdot \alpha \\
& \sim_{rat}
0 \quad (\text{par le cas initial})
\end{align*}
Cela achève la démonstration.
\end{proof}

\begin{remark}
\label{remark-quasi-projective}
Le Lemme \ref{lemma-moving-move} et le
Théorème \ref{theorem-well-defined}
valent aussi pour les variétés quasi-projectives non singulières, avec
la même démonstration. La seule modification consiste à démontrer la version suivante
du Lemme de déplacement \ref{lemma-moving} : Soit  $X \subset \mathbf{P}^N$
une sous-variété fermée. Soient $n = \dim(X)$ et $0 \leq d, d' < n$. Soit
$X^{reg} \subset X$ l'ouvert des points non singuliers. Soit
$Z \subset X^{reg}$ une sous-variété fermée de dimension $d$ et soient
$T_i \subset X^{reg}$, $i \in I$, une collection finie de sous-variétés fermées
de dimension $d'$. Alors il existe une sous-variété $C \subset \mathbf{P}^N$
telle que $C$ rencontre $X$ proprement et que
$$
(C \cdot X)|_{X^{reg}} = Z + \sum\nolimits_{j \in J} m_j Z_j
$$
où les $Z_j \subset X^{reg}$ sont irréductibles de dimension $d$, distinctes
de $Z$, et
$$
\dim(Z_j \cap T_i) \leq \dim(Z \cap T_i)
$$
avec inégalité stricte si $Z$ ne rencontre pas $T_i$ proprement dans $X^{reg}$.
\end{remark}



\section{Anneaux de Chow}
\label{section-chow-rings}

\noindent
Soit $X$ une variété projective non singulière. Nous définissons le produit
d'intersection
$$
\CH_r(X) \times \CH_s(X) \longrightarrow \CH_{r + s - \dim(X)}(X),\quad
(\alpha, \beta) \longmapsto \alpha \cdot \beta
$$
comme suit. Soient $\alpha \in Z_r(X)$ et $\beta \in Z_s(X)$.
Si $\alpha$ et $\beta$ se rencontrent proprement, nous utilisons la
définition donnée à la Section \ref{section-intersection-product}.
Sinon, choisissons $\alpha \sim_{rat} \alpha'$ comme dans le
Lemme \ref{lemma-moving-move} et posons
$$
\alpha \cdot \beta =
\text{classe de }\alpha' \cdot \beta \in \CH_{r + s - \dim(X)}(X)
$$
Ceci est bien défini et passe au quotient par l'équivalence rationnelle d'après le
Théorème \ref{theorem-well-defined}. Le produit d'intersection
sur $\CH_*(X)$ est commutatif (c'est clair), associatif
(Lemme \ref{lemma-associative}) et possède une unité $[X] \in \CH_{\dim(X)}(X)$.

\medskip\noindent
Nous employons souvent $\CH^c(X) = \CH_{\dim X - c}(X)$ pour désigner le groupe de Chow des
cycles de codimension $c$, voir
Homologie de Chow, Section \ref{chow-section-cycles-codimension}.
Le produit d'intersection définit un produit
$$
\CH^k(X) \times \CH^l(X) \longrightarrow \CH^{k + l}(X)
$$
qui est commutatif, associatif et possède une unité $1 = [X] \in \CH^0(X)$.


\section{Image réciproque par un morphisme quelconque}
\label{section-general-pullback}

\noindent
Soit $f : X \to Y$ un morphisme de variétés projectives non singulières.
Nous définissons
$$
f^* : \CH_k(Y) \to \CH_{k+\dim X - \dim Y}(X)
$$
par la règle
$$
f^*(\alpha) = pr_{X, *}(\Gamma_f \cdot pr_Y^*(\alpha))
$$
où $\Gamma_f \subset X\times Y$ est le graphe de $f$. Notons qu'à ce
degré de généralité, cette application n'est définie que sur les classes de cycles, et non sur les cycles. Avec la
notation $\CH^*$ introduite à la Section \ref{section-chow-rings},
on peut considérer l'image réciproque comme une application
$$
f^* : \CH^*(Y) \to \CH^*(X)
$$
autrement dit, comme une application de groupes abéliens gradués.

\begin{lemma}
\label{lemma-pullback-and-intersection-product}
Soit $f : X \to Y$ un morphisme de variétés projectives non singulières.
L'application d'image réciproque sur les groupes de Chow vérifie les propriétés suivantes :
\begin{enumerate}
\item $f^* : \CH^*(Y) \to \CH^*(X)$ est un morphisme d'anneaux,
\item $(g \circ f)^* = f^* \circ g^*$ pour un couple composable $f, g$,
\item la formule de projection vaut : $f_*(\alpha) \cdot \beta =
f_*( \alpha \cdot f^*\beta)$, et
\item si $f$ est plat, elle coïncide avec la définition précédente.
\end{enumerate}
\end{lemma}

\begin{proof}
Toutes ces assertions résultent immédiatement des résultats précédents.

\medskip\noindent
Pour (1), il suffit de montrer que
$\text{pr}_{X,*}( \Gamma_f \cdot \alpha \cdot \beta) =
\text{pr}_{X,*}(\Gamma_f \cdot \alpha) \cdot
\text{pr}_{X,*}(\Gamma_f \cdot \beta)$
pour des cycles $\alpha$, $\beta$ sur $X \times Y$. Si $\alpha$ est un cycle sur
$X \times Y$ qui rencontre $\Gamma_f$ proprement, alors il est facile
de voir que
$$
\Gamma_f \cdot \alpha =
\Gamma_f \cdot \text{pr}_X^*(\text{pr}_{X,*}(\Gamma_f \cdot \alpha))
$$
comme cycles, puisque $\Gamma_f$ est un graphe. On obtient ainsi la première
égalité dans
\begin{align*}
\text{pr}_{X,*}(\Gamma_f \cdot \alpha \cdot \beta)
& =
\text{pr}_{X,*}(
\Gamma_f \cdot
\text{pr}_X^*(\text{pr}_{X,*}(\Gamma_f \cdot \alpha)) \cdot \beta) \\
& =
\text{pr}_{X,*}(\text{pr}_X^*(\text{pr}_{X,*}(\Gamma_f \cdot \alpha))
\cdot (\Gamma_f \cdot \beta)) \\
& =
\text{pr}_{X,*}(\Gamma_f \cdot \alpha) \cdot
\text{pr}_{X,*}(\Gamma_f \cdot \beta)
\end{align*}
la dernière étape résultant de la formule de projection dans le cas plat
(Lemme \ref{lemma-projection-formula-flat}).

\medskip\noindent
Si $g : Y \to Z$, la propriété (2) résulte formellement de l'observation suivante :
$$
\Gamma =
\text{pr}_{X \times Y}^*\Gamma_f \cdot
\text{pr}_{Y \times Z}^*\Gamma_g
$$
dans $Z_*(X \times Y \times Z)$, où $\Gamma = \{(x, f(x), g(f(x))\}$,
et s'envoie isomorphiquement sur $\Gamma_{g \circ f}$ dans $X \times Z$.
L'égalité résulte de l'égalité schématique et du
Lemme \ref{lemma-transversal}.

\medskip\noindent
Pour (3), nous appliquons deux fois la formule de projection pour les applications plates :
\begin{align*}
f_*(\alpha \cdot pr_{X, *}(\Gamma_f \cdot pr_Y^*(\beta)))
& =
f_*(pr_{X, *}(pr_X^*\alpha \cdot \Gamma_f \cdot pr_Y^*(\beta))) \\
& =
pr_{Y, *}(pr_X^*\alpha \cdot \Gamma_f \cdot pr_Y^*(\beta))) \\
& =
pt_{Y, *}(pr_X^*\alpha \cdot \Gamma_f) \cdot \beta \\
& =
f_*(\alpha) \cdot \beta
\end{align*}
où, dans la dernière égalité, nous utilisons la remarque précédente sur les graphes.
Cela démontre (3).

\medskip\noindent
La propriété (4) repose sur l'identification du produit d'intersection
$\Gamma_f \cdot pr_Y^*\alpha$ lorsque $f$ est plat. Plus précisément, dans ce
cas, si $V \subset Y$ est une sous-variété fermée, tout point générique
$\xi$ du schéma $f^{-1}(V) \cong \Gamma_f \cap pr_Y^{-1}(V)$
est au-dessus du point générique de $V$. L'anneau local de
$pr_Y^{-1}(V) = X \times V$ en $\xi$ est donc de Cohen--Macaulay. Comme
$\Gamma_f \subset X \times Y$ est une immersion régulière (en tant que morphisme de
variétés projectives lisses), on obtient
$$
\Gamma_f \cdot pr_Y^*[V] = [\Gamma_f \cap pr_Y^{-1}(V)]_d
$$
où $d$ est la dimension de $\Gamma_f \cap pr_Y^{-1}(V)$, voir le
Lemme \ref{lemma-multiplicity-lci-CM}. Comme $\Gamma_f \cap pr_Y^{-1}(V)$
s'envoie isomorphiquement sur $f^{-1}(V)$, on conclut.
\end{proof}


\section{Image réciproque des cycles}
\label{section-pullback-cycles}

\noindent
Supposons que $X$ et $Y$ soient des variétés
projectives non singulières, et soit $f : X \to Y$ un morphisme.
Supposons que $Z \subset Y$ soit une sous-variété fermée. Soit $f^{-1}(Z)$
l'image réciproque schématique :
$$
\xymatrix{
f^{-1}(Z) \ar[r] \ar[d] & Z \ar[d] \\
X \ar[r] & Y
}
$$
est un diagramme cartésien de schémas. En particulier, $f^{-1}(Z) \subset X$
est un sous-schéma fermé de $X$. Dans ce cas, on a toujours
$$
\dim f^{-1}(Z) \geq \dim Z + \dim X - \dim Y.
$$
Si l'égalité a lieu dans la formule précédente, alors
\[f^*[Z] = [f^{-1}(Z)]_{\dim Z + \dim X - \dim Y}\]
pourvu que le schéma $Z$ soit de Cohen--Macaulay aux images
des points génériques de $f^{-1}(Z)$. Cela résulte de l'identification de
$f^{-1}(Z)$ avec l'intersection schématique de $\Gamma_f$
et $X \times Z$, puis du Lemme \ref{lemma-multiplicity-lci-CM}.
Les détails sont analogues à la démonstration de la partie (4) du
Lemme \ref{lemma-pullback-and-intersection-product} ci-dessus.





\begin{multicols}{2}[\section{Autres chapitres}]
\noindent
Préliminaires
\begin{enumerate}
\item \hyperref[introduction-section-phantom]{Introduction}
\item \hyperref[conventions-section-phantom]{Conventions}
\item \hyperref[sets-section-phantom]{Théorie des ensembles}
\item \hyperref[categories-section-phantom]{Catégories}
\item \hyperref[topology-section-phantom]{Topologie}
\item \hyperref[sheaves-section-phantom]{Faisceaux sur les espaces}
\item \hyperref[sites-section-phantom]{Sites et faisceaux}
\item \hyperref[stacks-section-phantom]{Champs}
\item \hyperref[fields-section-phantom]{Corps}
\item \hyperref[algebra-section-phantom]{Algèbre commutative}
\item \hyperref[brauer-section-phantom]{Groupes de Brauer}
\item \hyperref[homology-section-phantom]{Algèbre homologique}
\item \hyperref[derived-section-phantom]{Catégories dérivées}
\item \hyperref[simplicial-section-phantom]{Méthodes simpliciales}
\item \hyperref[more-algebra-section-phantom]{Compléments d'algèbre}
\item \hyperref[smoothing-section-phantom]{Lissage des morphismes d'anneaux}
\item \hyperref[modules-section-phantom]{Faisceaux de modules}
\item \hyperref[sites-modules-section-phantom]{Modules sur les sites}
\item \hyperref[injectives-section-phantom]{Objets injectifs}
\item \hyperref[cohomology-section-phantom]{Cohomologie des faisceaux}
\item \hyperref[sites-cohomology-section-phantom]{Cohomologie sur les sites}
\item \hyperref[dga-section-phantom]{Algèbre différentielle graduée}
\item \hyperref[dpa-section-phantom]{Algèbre à puissances divisées}
\item \hyperref[sdga-section-phantom]{Faisceaux différentiels gradués}
\item \hyperref[hypercovering-section-phantom]{Hyperrecouvrements}
\end{enumerate}
Schémas
\begin{enumerate}
\setcounter{enumi}{25}
\item \hyperref[schemes-section-phantom]{Schémas}
\item \hyperref[constructions-section-phantom]{Constructions de schémas}
\item \hyperref[properties-section-phantom]{Propriétés des schémas}
\item \hyperref[morphisms-section-phantom]{Morphismes de schémas}
\item \hyperref[coherent-section-phantom]{Cohomologie des schémas}
\item \hyperref[divisors-section-phantom]{Diviseurs}
\item \hyperref[limits-section-phantom]{Limites de schémas}
\item \hyperref[varieties-section-phantom]{Variétés}
\item \hyperref[topologies-section-phantom]{Topologies sur les schémas}
\item \hyperref[descent-section-phantom]{Descente}
\item \hyperref[perfect-section-phantom]{Catégories dérivées de schémas}
\item \hyperref[more-morphisms-section-phantom]{Compléments sur les morphismes}
\item \hyperref[flat-section-phantom]{Compléments sur la platitude}
\item \hyperref[groupoids-section-phantom]{Schémas en groupoïdes}
\item \hyperref[more-groupoids-section-phantom]{Compléments sur les schémas en groupoïdes}
\item \hyperref[etale-section-phantom]{Morphismes étales de schémas}
\end{enumerate}
Thèmes de théorie des schémas
\begin{enumerate}
\setcounter{enumi}{41}
\item \hyperref[chow-section-phantom]{Homologie de Chow}
\item \hyperref[intersection-section-phantom]{Théorie de l'intersection}
\item \hyperref[pic-section-phantom]{Schémas de Picard des courbes}
\item \hyperref[weil-section-phantom]{Théories cohomologiques de Weil}
\item \hyperref[adequate-section-phantom]{Modules adéquats}
\item \hyperref[dualizing-section-phantom]{Complexes dualisants}
\item \hyperref[duality-section-phantom]{Dualité pour les schémas}
\item \hyperref[discriminant-section-phantom]{Discriminants et différentes}
\item \hyperref[derham-section-phantom]{Cohomologie de de Rham}
\item \hyperref[local-cohomology-section-phantom]{Cohomologie locale}
\item \hyperref[algebraization-section-phantom]{Géométrie algébrique et formelle}
\item \hyperref[curves-section-phantom]{Courbes algébriques}
\item \hyperref[resolve-section-phantom]{Résolution des surfaces}
\item \hyperref[models-section-phantom]{Réduction semi-stable}
\item \hyperref[functors-section-phantom]{Foncteurs et morphismes}
\item \hyperref[equiv-section-phantom]{Catégories dérivées de variétés}
\item \hyperref[pione-section-phantom]{Groupes fondamentaux de schémas}
\item \hyperref[etale-cohomology-section-phantom]{Cohomologie étale}
\item \hyperref[crystalline-section-phantom]{Cohomologie cristalline}
\item \hyperref[proetale-section-phantom]{Cohomologie pro-étale}
\item \hyperref[relative-cycles-section-phantom]{Cycles relatifs}
\item \hyperref[more-etale-section-phantom]{Compléments de cohomologie étale}
\item \hyperref[trace-section-phantom]{La formule des traces}
\end{enumerate}
Espaces algébriques
\begin{enumerate}
\setcounter{enumi}{64}
\item \hyperref[spaces-section-phantom]{Espaces algébriques}
\item \hyperref[spaces-properties-section-phantom]{Propriétés des espaces algébriques}
\item \hyperref[spaces-morphisms-section-phantom]{Morphismes d'espaces algébriques}
\item \hyperref[decent-spaces-section-phantom]{Espaces algébriques décents}
\item \hyperref[spaces-cohomology-section-phantom]{Cohomologie des espaces algébriques}
\item \hyperref[spaces-limits-section-phantom]{Limites d'espaces algébriques}
\item \hyperref[spaces-divisors-section-phantom]{Diviseurs sur les espaces algébriques}
\item \hyperref[spaces-over-fields-section-phantom]{Espaces algébriques sur des corps}
\item \hyperref[spaces-topologies-section-phantom]{Topologies sur les espaces algébriques}
\item \hyperref[spaces-descent-section-phantom]{Descente et espaces algébriques}
\item \hyperref[spaces-perfect-section-phantom]{Catégories dérivées d'espaces}
\item \hyperref[spaces-more-morphisms-section-phantom]{Compléments sur les morphismes d'espaces}
\item \hyperref[spaces-flat-section-phantom]{Platitude sur les espaces algébriques}
\item \hyperref[spaces-groupoids-section-phantom]{Groupoïdes dans les espaces algébriques}
\item \hyperref[spaces-more-groupoids-section-phantom]{Compléments sur les groupoïdes dans les espaces}
\item \hyperref[bootstrap-section-phantom]{Amorçage}
\item \hyperref[spaces-pushouts-section-phantom]{Sommes amalgamées d'espaces algébriques}
\end{enumerate}
Thèmes de géométrie
\begin{enumerate}
\setcounter{enumi}{81}
\item \hyperref[spaces-chow-section-phantom]{Groupes de Chow des espaces}
\item \hyperref[groupoids-quotients-section-phantom]{Quotients de groupoïdes}
\item \hyperref[spaces-more-cohomology-section-phantom]{Compléments sur la cohomologie des espaces}
\item \hyperref[spaces-simplicial-section-phantom]{Espaces simpliciaux}
\item \hyperref[spaces-duality-section-phantom]{Dualité pour les espaces}
\item \hyperref[formal-spaces-section-phantom]{Espaces algébriques formels}
\item \hyperref[restricted-section-phantom]{Algébrisation des espaces formels}
\item \hyperref[spaces-resolve-section-phantom]{Résolution des surfaces revisitée}
\end{enumerate}
Théorie des déformations
\begin{enumerate}
\setcounter{enumi}{89}
\item \hyperref[formal-defos-section-phantom]{Théorie formelle des déformations}
\item \hyperref[defos-section-phantom]{Théorie des déformations}
\item \hyperref[cotangent-section-phantom]{Le complexe cotangent}
\item \hyperref[examples-defos-section-phantom]{Problèmes de déformation}
\end{enumerate}
Champs algébriques
\begin{enumerate}
\setcounter{enumi}{93}
\item \hyperref[algebraic-section-phantom]{Champs algébriques}
\item \hyperref[examples-stacks-section-phantom]{Exemples de champs}
\item \hyperref[stacks-sheaves-section-phantom]{Faisceaux sur les champs algébriques}
\item \hyperref[criteria-section-phantom]{Critères de représentabilité}
\item \hyperref[artin-section-phantom]{Axiomes d'Artin}
\item \hyperref[quot-section-phantom]{Espaces de Quot et de Hilbert}
\item \hyperref[stacks-properties-section-phantom]{Propriétés des champs algébriques}
\item \hyperref[stacks-morphisms-section-phantom]{Morphismes de champs algébriques}
\item \hyperref[stacks-limits-section-phantom]{Limites de champs algébriques}
\item \hyperref[stacks-cohomology-section-phantom]{Cohomologie des champs algébriques}
\item \hyperref[stacks-perfect-section-phantom]{Catégories dérivées de champs}
\item \hyperref[stacks-introduction-section-phantom]{Introduction aux champs algébriques}
\item \hyperref[stacks-more-morphisms-section-phantom]{Compléments sur les morphismes de champs}
\item \hyperref[stacks-geometry-section-phantom]{La géométrie des champs}
\end{enumerate}
Thèmes de théorie des espaces de modules
\begin{enumerate}
\setcounter{enumi}{107}
\item \hyperref[moduli-section-phantom]{Champs de modules}
\item \hyperref[moduli-curves-section-phantom]{Espaces de modules de courbes}
\end{enumerate}
Divers
\begin{enumerate}
\setcounter{enumi}{109}
\item \hyperref[examples-section-phantom]{Exemples}
\item \hyperref[exercises-section-phantom]{Exercices}
\item \hyperref[guide-section-phantom]{Guide bibliographique}
\item \hyperref[desirables-section-phantom]{Desiderata}
\item \hyperref[coding-section-phantom]{Style de programmation}
\item \hyperref[obsolete-section-phantom]{Obsolète}
\item \hyperref[fdl-section-phantom]{Licence de documentation libre GNU}
\item \hyperref[index-section-phantom]{Index généré automatiquement}
\end{enumerate}
\end{multicols}


\bibliography{my}
\bibliographystyle{amsalpha}

\end{document}
